% chapters/ch03.tex
% 第~
\ref{chap:3}~章 \index{ruotuopu@弱拓扑 (weak topology)}
弱拓扑. \index{zifankongjian@自反空间 (reflexive space)}
自反空间. \index{kefenkongjian@可分空间 (separable space)}
可分空间. \index{yizhitukongjian@一致凸空间 (uniformly convex space)}
一致凸性
% 源: 原书 PDF 第 65--97 页 (印刷页 55--87)

\chapter{弱拓扑. 自反空间. 可分空间. 一致凸性}
\label{chap:3}

\section{使一族映射连续的, 最粗的拓扑}
\label{sec:3.1}

本章开头先回顾拓扑学中一个熟知的概念.设 $X$ 为一个集合 (不带任何结构),
$(Y_i)_{i\in I}$ 为一族\emph{拓扑空间}.给定一族映射 $(\varphi_i)_{i\in I}$, 其中对
每个 $i\in I$, $\varphi_i$ 都把 $X$ 映入 $Y_i$.我们考虑下述问题:

\runin{问题 1.}
在 $X$ 上构造一个拓扑, 使所有映射 $(\varphi_i)_{i\in I}$ 都连续.如果可能,
找出一个\emph{最经济}的拓扑 $\mathcal{T}$, 即开集最少的拓扑.

注意: 若给 $X$ 配上离散拓扑 (即 $X$ 的每个子集都是开集), 则每个映射
$\varphi_i$ 都连续; 当然, 这个拓扑远非``最便宜''的, 事实上它是最贵的! 我们将会
看到, 总存在一个 (唯一的) ``最便宜''的拓扑 $\mathcal{T}$, 使每个映射 $\varphi_i$ 都
连续.它称为与族 $(\varphi_i)_{i\in I}$ 相伴的\emph{最粗}或\emph{最弱}拓扑 (有时
也称为\emph{初始拓扑}).

若 $\omega_i\subset Y_i$ 是任一开集, 则 $\varphi_i^{-1}(\omega_i)$ 必须是 $\mathcal{T}$
中的开集.当 $\omega_i$ 取遍 $Y_i$ 的开集族, $i$ 取遍 $I$ 时, 我们得到一族 $X$ 的
子集, 其中每一个在拓扑 $\mathcal{T}$ 下都必须是开集.把这族子集记为
$(U_\lambda)_{\lambda\in\Lambda}$.当然, 这族子集不必构成一个拓扑.因此我们引出:

\runin{问题 2.}
给定集合 $X$ 以及 $X$ 中的一族子集 $(U_\lambda)_{\lambda\in\Lambda}$, 在 $X$ 上构造
最便宜的拓扑 $\mathcal{T}$, 使 $U_\lambda$ 对一切 $\lambda\in\Lambda$ 都是开集.

换言之, 我们要找出 $X$ 的子集构成的最便宜的一族 $\mathcal{F}$, 它对有限交与任意并
稳定\footnote{意思是: $\mathcal{F}$ 中有限多个集合的交以及 $\mathcal{F}$ 中任意多个
集合的并都仍属于 $\mathcal{F}$.}, 并且具有性质: 对每个 $\lambda\in\Lambda$ 都有
$U_\lambda\in\mathcal{F}$.构造如下.首先考虑 $(U_\lambda)_{\lambda\in\Lambda}$ 中
集合的有限交, 即 $\bigcap_{\lambda\in\Gamma}U_\lambda$, 其中 $\Gamma\subset\Lambda$
为有限集.这样得到 $X$ 的一族新子集, 记作 $\Phi$, 它包含
$(U_\lambda)_{\lambda\in\Lambda}$, 并且对有限交稳定.然而它不必对任意并稳定.因此
我们接着考虑由 $\Phi$ 中元素作任意并所得的一族 $\mathcal{F}$.显然 $\mathcal{F}$
对任意并稳定.$\mathcal{F}$ 是否对有限交稳定则并不显然; 但事实上我们有下述结果:

\begin{lemma}
\label{lem:3.1}
族 $\mathcal{F}$ 对有限交稳定.
\end{lemma}

引理  \ref{lem:3.1} 的证明——一道令人愉快的集合论练习——留给读者; 参见例如
G.~Folland \bref{Folland}{2}.现在显然, 上述构造给出了具有所需性质的最便宜的
拓扑.

\begin{remark}
\label{rem:3.1}
在 $\mathcal{F}$ 的构造中不能交换运算的次序.同样自然的是先取任意并, 再取有限交.
这样得到的一族对有限交稳定; 但它对任意并并不稳定.必须再考虑一次任意并, 过程
才稳定下来.
\end{remark}

总结上述讨论可知: 拓扑 $\mathcal{T}$ 的开集, 是先取形如
$\varphi_i^{-1}(\omega_i)$ 的集合的有限交, 再取任意并而得到的.由此可知, 对每个
$x\in X$, 取形如 $\bigcap_{\text{有限}}\varphi_i^{-1}(V_i)$ 的集合 (其中 $V_i$ 是
$\varphi_i(x)$ 在 $Y_i$ 中的邻域) 便得到 $x$ 在拓扑 $\mathcal{T}$ 下的一个邻域基.
回忆: 在拓扑空间中, 点 $x$ 的\emph{邻域基}是指 $x$ 的一族邻域, 使得 $x$ 的每个
邻域都包含该族中的某个邻域.

以下总给 $X$ 配上与族 $(\varphi_i)_{i\in I}$ 相伴的最弱拓扑 $\mathcal{T}$.拓扑
$\mathcal{T}$ 有两条简单的性质:

\begin{keyproposition}
\label{prop:3.1}
设 $(x_n)$ 为 $X$ 中的序列.则 $x_n\to x$ (在 $\mathcal{T}$ 中) 当且仅当对每个
$i\in I$ 都有 $\varphi_i(x_n)\to\varphi_i(x)$.
\end{keyproposition}

\begin{proof}
若 $x_n\to x$, 则对每个 $i$ 都有 $\varphi_i(x_n)\to\varphi_i(x)$, 因为每个
$\varphi_i$ 关于 $\mathcal{T}$ 连续.反之, 设 $U$ 为 $x$ 的邻域.由前面的讨论, 总可
设 $U$ 具有形式 $U=\bigcap_{i\in J}\varphi_i^{-1}(V_i)$, 其中 $J\subset I$ 为有限集.
对每个 $i\in J$ 存在整数 $N_i$, 使得当 $n\ge N_i$ 时 $\varphi_i(x_n)\in V_i$.
于是当 $n\ge N=\max_{i\in J}N_i$ 时 $x_n\in U$.
\end{proof}

\begin{keyproposition}
\label{prop:3.2}
设 $Z$ 为拓扑空间, $\psi$ 为从 $Z$ 到 $X$ 的映射.则 $\psi$ 连续当且仅当对每个
$i\in I$, $\varphi_i\circ\psi$ 都是从 $Z$ 到 $Y_i$ 的连续映射.
\end{keyproposition}

\begin{proof}
若 $\psi$ 连续, 则对每个 $i\in I$, $\varphi_i\circ\psi$ 也连续.反之, 须证明对每个
开集 $U\subset X$, $\psi^{-1}(U)$ 在 $Z$ 中开.但已知 $U$ 具有形式
$U=\bigcup_{\text{任意}}\bigcap_{\text{有限}}\varphi_i^{-1}(\omega_i)$, 其中
$\omega_i$ 在 $Y_i$ 中开.于是
\begin{equation*}
\psi^{-1}(U)
=\bigcup_{\text{任意}}\bigcap_{\text{有限}}\psi^{-1}\bigl[\varphi_i^{-1}(\omega_i)\bigr]
=\bigcup_{\text{任意}}\bigcap_{\text{有限}}(\varphi_i\circ\psi)^{-1}(\omega_i),
\end{equation*}
它在 $Z$ 中开, 因为每个映射 $\varphi_i\circ\psi$ 都连续.
\end{proof}

\section{弱拓扑 $\sigma(E,E^\star)$ 的定义与初等性质}
\label{sec:3.2}

设 $E$ 为 Banach 空间, $f\in E^\star$.以 $\varphi_f:E\to\R$ 记线性泛函
$\varphi_f(x)=\dual{f}{x}$.当 $f$ 取遍 $E^\star$ 时, 我们得到一族从 $E$ 到 $\R$
的映射 $(\varphi_f)_{f\in E^\star}$.现在暂且不管 $E$ 上通常的拓扑 (与 $\norm{\cdot}$
相联系的拓扑), 而在集合 $E$ 上定义一个新拓扑如下:

\begin{definition}
$E$ 上的\emph{弱拓扑} $\sigma(E,E^\star)$ 是与族 $(\varphi_f)_{f\in E^\star}$
相伴的最粗拓扑 (按第  \ref{sec:3.1} 节的意义, 其中 $X=E$, 每个 $Y_i=\R$, 且
$I=E^\star$).
\end{definition}

注意每个映射 $\varphi_f$ 关于通常拓扑连续, 因此弱拓扑比通常拓扑更弱.

\begin{proposition}
\label{prop:3.3}
弱拓扑 $\sigma(E,E^\star)$ 是 Hausdorff 拓扑.
\end{proposition}

\begin{proof}
给定 $x_1,x_2\in E$, $x_1\ne x_2$, 须找出弱拓扑 $\sigma(E,E^\star)$ 中的两个开集
$O_1$ 与 $O_2$, 使 $x_1\in O_1$, $x_2\in O_2$, 且 $O_1\cap O_2=\emptyset$.
由 Hahn--Banach 定理 (第二几何形式), 存在一个闭超平面把 $\{x_1\}$ 与 $\{x_2\}$
严格分离.于是存在 $f\in E^\star$ 与 $\alpha\in\R$, 使
\begin{equation*}
\dual{f}{x_1}<\alpha<\dual{f}{x_2}.
\end{equation*}
置
\begin{equation*}
O_1=\{x\in E;\ \dual{f}{x}<\alpha\}=\varphi_f^{-1}\bigl((-\infty,\alpha)\bigr),
\end{equation*}
\begin{equation*}
O_2=\{x\in E;\ \dual{f}{x}>\alpha\}=\varphi_f^{-1}\bigl((\alpha,+\infty)\bigr).
\end{equation*}
显然 $O_1$ 与 $O_2$ 关于 $\sigma(E,E^\star)$ 开, 且满足所需性质.
\end{proof}

\begin{keyproposition}
\label{prop:3.4}
设 $x_0\in E$; 给定 $\varepsilon>0$ 与 $E^\star$ 中的\emph{有限}集
$\{f_1,f_2,\dots,f_k\}$, 考虑
\begin{equation*}
V=V(f_1,f_2,\dots,f_k;\varepsilon)
=\{x\in E;\ \abs{\dual{f_i}{x-x_0}}<\varepsilon\ \ \forall i=1,2,\dots,k\}.
\end{equation*}
则 $V$ 是 $x_0$ 在拓扑 $\sigma(E,E^\star)$ 下的邻域.此外, 令 $\varepsilon$, $k$ 与
$E^\star$ 中的 $f_i$ 变化, 便得到 $x_0$ 关于 $\sigma(E,E^\star)$ 的一个邻域基.
\end{keyproposition}

\begin{proof}
显然 $V=\bigcap_{i=1}^k\varphi_{f_i}^{-1}\bigl((a_i-\varepsilon,a_i+\varepsilon)\bigr)$,
其中 $a_i=\dual{f_i}{x_0}$, 它关于拓扑 $\sigma(E,E^\star)$ 开且包含 $x_0$.反之,
设 $U$ 为 $x_0$ 关于 $\sigma(E,E^\star)$ 的邻域.由第  \ref{sec:3.1} 节的讨论知,
存在包含 $x_0$ 的开集 $W$, $W\subset U$, 具有形式
$W=\bigcap_{\text{有限}}\varphi_{f_i}^{-1}(\omega_i)$, 其中 $\omega_i$ 是
$a_i=\dual{f_i}{x_0}$ 在 $\R$ 中的邻域.于是存在 $\varepsilon>0$, 使对每个 $i$ 都有
$(a_i-\varepsilon,a_i+\varepsilon)\subset\omega_i$.由此 $x_0\in V\subset W\subset U$.
\end{proof}

\begin{notation}
若 $E$ 中的序列 $(x_n)$ 在弱拓扑 $\sigma(E,E^\star)$ 下收敛于 $x$, 我们写
\begin{equation*}
\boxed{x_n\weakto x.}
\end{equation*}
为避免混淆, 有时也说``$x_n\weakto x$ 在 $\sigma(E,E^\star)$ 中弱收敛''.为了完全
清楚, 有时也用``$x_n\to x$ 强收敛''来强调强收敛, 意即 $\norm{x_n-x}\to0$.
\end{notation}

\begin{keyproposition}
\label{prop:3.5}
设 $(x_n)$ 为 $E$ 中的序列.则
\begin{enumerate}[label=(\roman*), leftmargin=2.4em, itemsep=0.15em]
	\item $[x_n\weakto x\ \text{在}\ \sigma(E,E^\star)\ \text{中弱收敛}]
	      \iff[\dual{f}{x_n}\to\dual{f}{x}\ \forall f\in E^\star]$.
	\item 若 $x_n\to x$ 强收敛, 则 $x_n\weakto x$ 在 $\sigma(E,E^\star)$ 中弱收敛.
	\item 若 $x_n\weakto x$ 在 $\sigma(E,E^\star)$ 中弱收敛, 则 $(\norm{x_n})$ 有界,
	      且 $\norm{x}\le\liminf\norm{x_n}$.
	\item 若 $x_n\weakto x$ 在 $\sigma(E,E^\star)$ 中弱收敛, 且 $f_n\to f$ 在
	      $E^\star$ 中强收敛 (即 $\norm{f_n-f}_{E^\star}\to0$), 则
	      $\dual{f_n}{x_n}\to\dual{f}{x}$.
\end{enumerate}
\end{keyproposition}

\begin{proof}
(i) 由命题  \ref{prop:3.1} 与弱拓扑 $\sigma(E,E^\star)$ 的定义即得.

(ii) 由 (i) 即得, 因为
$\abs{\dual{f}{x_n}-\dual{f}{x}}\le\norm{f}\norm{x_n-x}$; 由弱拓扑比强拓扑更弱这一
事实也可看出.

(iii) 由一致有界原理 (见推论  \ref{cor:2.4}) 即得, 因为对每个 $f\in E^\star$, 集合
$(\dual{f}{x_n})_n$ 有界.在不等式
\begin{equation*}
\abs{\dual{f}{x_n}}\le\norm{f}\norm{x_n}
\end{equation*}
中取极限, 得
\begin{equation*}
\abs{\dual{f}{x}}\le\norm{f}\liminf\norm{x_n},
\end{equation*}
从而 (由推论  \ref{cor:1.4})
\begin{equation*}
\norm{x}=\sup_{\norm{f}\le1}\abs{\dual{f}{x}}\le\liminf\norm{x_n}.
\end{equation*}

(iv) 由不等式
\begin{equation*}
\abs{\dual{f_n}{x_n}-\dual{f}{x}}
\le\abs{\dual{f_n-f}{x_n}}+\abs{\dual{f}{x_n-x}}
\le\norm{f_n-f}\norm{x_n}+\abs{\dual{f}{x_n-x}},
\end{equation*}
结合 (i) 与 (iii) 即得.
\end{proof}

\begin{keyproposition}
\label{prop:3.6}
当 $E$ 为有限维时, 弱拓扑 $\sigma(E,E^\star)$ 与通常拓扑相同.特别地, 序列
$(x_n)$ 弱收敛当且仅当它强收敛.
\end{keyproposition}

\begin{proof}
由于弱拓扑的开集总比强拓扑少, 只须验证每个强开集都是弱开集.设 $x_0\in E$,
$U$ 为 $x_0$ 在强拓扑中的邻域.须找出 $x_0$ 在弱拓扑 $\sigma(E,E^\star)$ 中的
邻域 $V$, 使 $V\subset U$.换言之, 须找出 $E^\star$ 中的
$f_1,f_2,\dots,f_k$ 与 $\varepsilon>0$, 使
\begin{equation*}
V=\{x\in E;\ \abs{\dual{f_i}{x-x_0}}<\varepsilon\ \ \forall i=1,2,\dots,k\}\subset U.
\end{equation*}
取 $r>0$ 使 $B(x_0,r)\subset U$.在 $E$ 中取基 $e_1,e_2,\dots,e_k$, 使
$\norm{e_i}=1$, $\forall i$.每个 $x\in E$ 都有分解
$x=\sum_{i=1}^k x_ie_i$, 而映射 $x\mapsto x_i$ 是 $E$ 上的连续线性泛函, 记为
$f_i$.对每个 $x\in V$ 有
\begin{equation*}
\norm{x-x_0}\le\sum_{i=1}^k\abs{\dual{f_i}{x-x_0}}<k\varepsilon.
\end{equation*}
取 $\varepsilon=r/k$, 即得 $V\subset U$.
\end{proof}

\begin{remark}
\label{rem:3.2}
弱拓扑 $\sigma(E,E^\star)$ 中的开集 (相应地, 闭集) 总是强拓扑中的开集 (相应地,
闭集).在任何无穷维空间中, 弱拓扑都严格粗于强拓扑; 即存在强拓扑中的开集
(相应地, 闭集), 它们在弱拓扑中不是开集 (相应地, 闭集).下面给出两个例子:
\end{remark}

\begin{example}
\label{exa:3.1}
单位球面 $S=\{x\in E;\ \norm{x}=1\}$ (其中 $E$ 为无穷维) 在弱拓扑
$\sigma(E,E^\star)$ 中绝不闭.更确切地说, 有
\begin{equation}
\overline{S}^{\sigma(E,E^\star)}=B_E,
\label{eq:3.1}
\end{equation}
其中 $\overline{S}^{\sigma(E,E^\star)}$ 表示 $S$ 在拓扑 $\sigma(E,E^\star)$ 中的闭包,
而 $B_E$ (第~
\ref{chap:2}~章中已定义) 表示 $E$ 中的闭单位球,
\begin{equation*}
B_E=\{x\in E;\ \norm{x}\le1\}.
\end{equation*}

先验证每个满足 $\norm{x_0}<1$ 的 $x_0\in E$ 都属于
$\overline{S}^{\sigma(E,E^\star)}$.事实上, 设 $V$ 为 $x_0$ 在
$\sigma(E,E^\star)$ 中的邻域.须证明 $V\cap S\ne\emptyset$.由命题  \ref{prop:3.4},
总可设 $V$ 具有形式
\begin{equation*}
V=\{x\in E;\ \abs{\dual{f_i}{x-x_0}}<\varepsilon\ \ \forall i=1,2,\dots,k\},
\end{equation*}
其中 $\varepsilon>0$, $f_1,f_2,\dots,f_k\in E^\star$.取 $y_0\in E$, $y_0\ne0$, 使
\begin{equation*}
\dual{f_i}{y_0}=0\quad\forall i=1,2,\dots,k.
\end{equation*}
[这样的 $y_0$ 存在; 否则映射 $\varphi:E\to\R^k$,
$\varphi(x)=(\dual{f_i}{x})_{1\le i\le k}$ 将是单射, 而 $\varphi$ 将是 $E$ 到
$\varphi(E)$ 的同构, 从而 $\dim E\le k$, 与 $E$ 为无穷维的假设矛盾.]\footnote{这一
构造的几何解释如下: 当 $E$ 为无穷维时, $x_0$ 在拓扑 $\sigma(E,E^\star)$ 中的每个
邻域 $V$ 都包含一条过 $x_0$ 的直线, 甚至包含一个过 $x_0$ 的``巨大''仿射空间.}
函数 $g(t)=\norm{x_0+ty_0}$ 在 $[0,\infty)$ 上连续, 且 $g(0)<1$,
$\lim_{t\to+\infty}g(t)=+\infty$.因此存在 $t_0>0$, 使 $\norm{x_0+t_0y_0}=1$.
于是 $x_0+t_0y_0\in V\cap S$, 从而得到
\begin{equation*}
S\subset B_E\subset\overline{S}^{\sigma(E,E^\star)}.
\end{equation*}

为完成 \eqref{eq:3.1} 的证明, 只须知道 $B_E$ 在拓扑 $\sigma(E,E^\star)$ 中闭.
但
\begin{equation*}
B_E=\bigcap_{\substack{f\in E^\star\\ \norm{f}\le1}}\{x\in E;\ \abs{\dual{f}{x}}\le1\},
\end{equation*}
它是弱闭集的交.
\end{example}

\begin{example}
\label{exa:3.2}
单位球 $U=\{x\in E;\ \norm{x}<1\}$ (其中 $E$ 为无穷维) 在弱拓扑
$\sigma(E,E^\star)$ 中绝不\emph{开}.反设 $U$ 弱开.则其补集
$U^c=\{x\in E;\ \norm{x}\ge1\}$ 弱闭.于是 $S=B_E\cap U^c$ 也弱闭; 这与例
\ref{exa:3.1} 矛盾.
\end{example}

\begin{starremark}
\label{rem:3.3}
在无穷维空间中, 弱拓扑绝不可度量化, 即 $E$ 上不存在度量 (从而也不存在范数) 在
$E$ 上诱导出弱拓扑 $\sigma(E,E^\star)$; 见习题  \ref{exer:3.8}.然而, 我们后面
将会看到 (定理  \ref{thm:3.29}), 若 $E^\star$ 可分, 则可在 $E$ 上定义一个范数,
它在 $E$ 的\emph{有界}集上诱导出弱拓扑 $\sigma(E,E^\star)$.
\end{starremark}

\begin{starremark}
\label{rem:3.4}
通常, 在无穷维空间中存在弱收敛而不强收敛的序列.例如, 若 $E^\star$ 可分或 $E$
自反, 则可构造 $E$ 中的序列 $(x_n)$, 使 $\norm{x_n}=1$ 且 $x_n\weakto0$ 弱收敛
(见习题  \ref{exer:3.22}).然而, 也存在无穷维空间, 使每个弱收敛序列都强收敛.
例如 $\ell^1$ 就有这个不寻常的性质 (见问题 8).这样的空间相当``稀少'', 且有些
``病态''.这一奇怪的事实与注记 \ref{rem:3.2} 并不矛盾; 注记 \ref{rem:3.2} 断言:
在无穷维空间中, 弱拓扑与强拓扑总是不同的, 弱拓扑严格粗于强拓扑.请记住: 两个
\emph{度量} (或可度量化) 空间若有相同的收敛序列, 则它们的拓扑相同; 然而, 两个
\emph{拓扑}空间若有相同的收敛序列, 它们的拓扑却不必相同.
\end{starremark}

\section{弱拓扑, 凸集与线性算子}
\label{sec:3.3}

每个弱闭集都是强闭集, 而其逆在无穷维空间中不成立 (见注记 \ref{rem:3.2}).然而,
知道下述事实非常有用: 对凸集而言, 弱闭 $=$ 强闭:

\begin{keytheorem}
\label{thm:3.7}
设 $C$ 为 $E$ 的凸子集.则 $C$ 在弱拓扑 $\sigma(E,E^\star)$ 中闭当且仅当它在强
拓扑中闭.
\end{keytheorem}

\begin{proof}
设 $C$ 强闭, 要证 $C$ 弱闭.我们将验证 $C$ 的补集 $C^c$ 在弱拓扑中开.为此,
设 $x_0\notin C$.由 Hahn--Banach 定理, 存在一个闭超平面把 $\{x_0\}$ 与 $C$
严格分离.于是存在 $f\in E^\star$ 与 $\alpha\in\R$, 使
\begin{equation*}
\dual{f}{x_0}<\alpha<\dual{f}{y}\quad\forall y\in C.
\end{equation*}
置
\begin{equation*}
V=\{x\in E;\ \dual{f}{x}<\alpha\};
\end{equation*}
于是 $x_0\in V$, $V\cap C=\emptyset$ (即 $V\subset C^c$), 且 $V$ 在弱拓扑中开.
\end{proof}

\begin{corollary}[Mazur]
\label{cor:3.8}
假设 $(x_n)$ 弱收敛于 $x$.则存在由 $x_n$ 的凸组合构成的序列 $(y_n)$, 它强收敛
于 $x$.
\end{corollary}

\begin{proof}
置 $C=\operatorname{conv}\bigl(\bigcup_{p=1}^\infty\{x_p\}\bigr)$, 即 $x_n$ 的凸包.
由于 $x$ 属于 $\bigcup_{p=1}^\infty\{x_p\}$ 的弱闭包, 它自然属于 $C$ 的弱闭包.
由定理  \ref{thm:3.7}, $x\in\overline{C}$ (即 $C$ 的强闭包), 结论由此得出.
\end{proof}

\begin{remark}
\label{rem:3.5}
推论  \ref{cor:3.8} 有一些变形 (见习题  \ref{exer:3.4} 与  \ref{exer:5.24}).此外,
注意定理  \ref{thm:3.7} 的证明表明: 每个闭凸集 $C$ 都等于所有包含 $C$ 的闭半空间
之交.
\end{remark}

\begin{keycorollary}
\label{cor:3.9}
假设 $\varphi:E\to(-\infty,+\infty]$ 凸, 且在强拓扑中下半连续.则 $\varphi$ 在弱
拓扑 $\sigma(E,E^\star)$ 中下半连续.
\end{keycorollary}

\begin{proof}
对每个 $\lambda\in\R$, 集合
\begin{equation*}
A=\{x\in E;\ \varphi(x)\le\lambda\}
\end{equation*}
凸且强闭.由定理  \ref{thm:3.7} 它弱闭, 因此 $\varphi$ 弱下半连续.
\end{proof}

\begin{keyremark}
\label{rem:3.6}
在实践中要证明一个函数在弱拓扑中下半连续可能相当困难.推论  \ref{cor:3.9} 常按
下述方式使用:
\begin{equation*}
\boxed{\varphi\ \text{凸且强连续}\ \Rightarrow\ \varphi\ \text{弱下半连续}.}
\end{equation*}
例如, 函数 $\varphi(x)=\norm{x}$ 凸且强连续; 因此它弱下半连续.特别地, 若
$x_n\weakto x$ 弱收敛, 则推出 $\norm{x}\le\liminf\norm{x_n}$ (另见命题
 \ref{prop:3.5}).
\end{keyremark}

\begin{theorem}
\label{thm:3.10}
设 $E$ 与 $F$ 为两个 Banach 空间, $T$ 为从 $E$ 到 $F$ 的线性算子.假设 $T$ 在强
拓扑中连续.则 $T$ 是从 $E$ 弱 $\sigma(E,E^\star)$ 到 $F$ 弱 $\sigma(F,F^\star)$
的连续算子, 反之亦然.
\end{theorem}

\begin{proof}
由命题  \ref{prop:3.2}, 只须验证对每个 $f\in F^\star$, 映射
$x\mapsto\dual{f}{Tx}$ 是从 $E$ 弱 $\sigma(E,E^\star)$ 到 $\R$ 的连续映射.但
$x\mapsto\dual{f}{Tx}$ 是 $E$ 上的连续线性泛函, 因此它在弱拓扑
$\sigma(E,E^\star)$ 中也连续.

反之, 设 $T$ 是从 $E$ 弱到 $F$ 弱的连续算子.则 $G(T)$ 在赋以乘积拓扑
$\sigma(E,E^\star)\times\sigma(F,F^\star)$ 的 $E\times F$ 中闭, 而该乘积拓扑显然
就是 $\sigma(E\times F,(E\times F)^\star)$.于是 $G(T)$ 强闭 (任何弱闭集都强闭).
借助闭图像定理 (定理  \ref{thm:2.9}), 我们断定 $T$ 是从 $E$ 强到 $F$ 强的连续
算子.
\end{proof}

\begin{remark}
\label{rem:3.7}
上述论证还表明: 若线性算子 $T$ 是从 $E$ 强到 $F$ 弱的连续算子, 则 $T$ 是从 $E$
强到 $F$ 强的连续算子.因此, 对线性算子而言, 下述连续性都是同一回事:
$S\to S$, $W\to W$, $S\to W$ ($S=$ 强, $W=$ 弱).另一方面, 很少有线性算子是
$W\to S$ 连续的; 这种情况成立当且仅当 $T$ 是 $S\to S$ 连续的, 并且
$\dim R(T)<\infty$ (见习题  \ref{exer:6.7}).

此外, 注意一般说来, 从 $E$ 强到 $F$ 强连续的非线性映射, 不必是从 $E$ 弱到 $F$ 弱
连续的 (见例如习题  \ref{exer:4.20}).这是非线性问题中困难的一个主要来源.
\end{remark}

\section{\texorpdfstring{\wstar{}拓扑 $\sigma(E^\star,E)$}{弱*拓扑 $\sigma(E^\star,E)$}}
\label{sec:3.4}

到目前为止, $E^\star$ 上有两个拓扑:
\begin{enumerate}[label=(\alph*), leftmargin=2.0em, itemsep=0.15em]
	\item 与 $E^\star$ 的范数相联系的通常 (强) 拓扑;
	\item 在 $E^\star$ 上施行第  \ref{sec:3.3} 节的构造所得的弱拓扑
	      $\sigma(E^\star,E^{\star\star})$.
\end{enumerate}
现在要在 $E^\star$ 上定义第三个拓扑, 称为\emph{\wstar{}拓扑}, 记作
$\sigma(E^\star,E)$ (这里的 $^\star$ 提醒我们: 这个拓扑只在共轭空间上定义).
对每个 $x\in E$ 考虑线性泛函 $\varphi_x:E^\star\to\R$,
$f\mapsto\varphi_x(f)=\dual{f}{x}$.当 $x$ 取遍 $E$ 时, 我们得到一族从 $E^\star$
到 $\R$ 的映射 $(\varphi_x)_{x\in E}$.

\begin{definition}
\emph{\wstar{}拓扑} $\sigma(E^\star,E)$ 是 $E^\star$ 上与族 $(\varphi_x)_{x\in E}$
相伴的最粗拓扑 (按第  \ref{sec:3.1} 节的意义, 其中 $X=E^\star$, 所有 $Y_i=\R$,
且 $I=E$).
\end{definition}

由于 $E\subset E^{\star\star}$, 显然拓扑 $\sigma(E^\star,E)$ 比拓扑
$\sigma(E^\star,E^{\star\star})$ 更粗; 即拓扑 $\sigma(E^\star,E)$ 的开集
(相应地, 闭集) 比拓扑 $\sigma(E^\star,E^{\star\star})$ 少, 而后者又比强拓扑的开集
(相应地, 闭集) 少.

\begin{remark}
\label{rem:3.8}
读者大概会问, 为什么人们对弱拓扑如此狂热! 原因如下: 更粗的拓扑有更多的紧集.
例如 $E^\star$ 中的闭单位球 $B_{E^\star}$ 在强拓扑中绝不紧 (除非
$\dim E<\infty$; 见定理  \ref{thm:6.5}), 但它在\wstar{}拓扑中总是紧的 (见定理
 \ref{thm:3.16}).了解了紧集的基本作用——例如在极小化这样的存在性机制中的作用——
就很容易理解\wstar{}拓扑的重要性.
\end{remark}

\begin{proposition}
\label{prop:3.11}
\wstar{}拓扑是 Hausdorff 拓扑.
\end{proposition}

\begin{proof}
给定 $f_1,f_2\in E^\star$, $f_1\ne f_2$, 存在 $x\in E$ 使
$\dual{f_1}{x}\ne\dual{f_2}{x}$ (这里不用 Hahn--Banach 定理, 只用到 $f_1\ne f_2$
这一事实).例如设 $\dual{f_1}{x}<\dual{f_2}{x}$, 取 $\alpha$ 使
\begin{equation*}
\dual{f_1}{x}<\alpha<\dual{f_2}{x}.
\end{equation*}
置
\begin{equation*}
O_1=\{f\in E^\star;\ \dual{f}{x}<\alpha\}=\varphi_x^{-1}\bigl((-\infty,\alpha)\bigr),
\end{equation*}
\begin{equation*}
O_2=\{f\in E^\star;\ \dual{f}{x}>\alpha\}=\varphi_x^{-1}\bigl((\alpha,+\infty)\bigr).
\end{equation*}
则 $O_1$ 与 $O_2$ 是 $\sigma(E^\star,E)$ 中的开集, 且 $f_1\in O_1$, $f_2\in O_2$,
$O_1\cap O_2=\emptyset$.
\end{proof}

\begin{proposition}
\label{prop:3.12}
设 $f_0\in E^\star$; 给定 $E$ 中的\emph{有限}集 $\{x_1,x_2,\dots,x_k\}$ 与
$\varepsilon>0$, 考虑
\begin{equation*}
V=V(x_1,x_2,\dots,x_k;\varepsilon)
=\bigl\{f\in E^\star;\ \abs{\dual{f-f_0}{x_i}}<\varepsilon\ \ \forall i=1,2,\dots,k\bigr\}.
\end{equation*}
则 $V$ 是 $f_0$ 在拓扑 $\sigma(E^\star,E)$ 下的邻域.此外, 令 $\varepsilon$, $k$ 与
$E$ 中的 $x_i$ 变化, 便得到 $f_0$ 关于 $\sigma(E^\star,E)$ 的一个邻域基.
\end{proposition}

\begin{proof}
与命题  \ref{prop:3.4} 的证明相同.
\end{proof}

\begin{notation}
若 $E^\star$ 中的序列 $(f_n)$ 在\wstar{}拓扑下收敛于 $f$, 我们写
\begin{equation*}
\boxed{f_n\weakstar f.}
\end{equation*}
为避免混淆, 有时也说``$f_n\weakstar f$ 在 $\sigma(E^\star,E)$ 中'', ``$f_n\weakto f$ 在 $\sigma(E^\star,E^{\star\star})$ 中''以及``$f_n\to f$ 强收敛''.
\end{notation}

\begin{keyproposition}
\label{prop:3.13}
设 $(f_n)$ 为 $E^\star$ 中的序列.则
\begin{enumerate}[label=(\roman*), leftmargin=2.4em, itemsep=0.15em]
	\item $[f_n\weakstar f\ \text{在}\ \sigma(E^\star,E)\ \text{中}]
	      \iff[\dual{f_n}{x}\to\dual{f}{x},\ \forall x\in E]$.
	\item 若 $f_n\to f$ 强收敛, 则 $f_n\weakto f$ 在
	      $\sigma(E^\star,E^{\star\star})$ 中.

	      若 $f_n\weakto f$ 在 $\sigma(E^\star,E^{\star\star})$ 中, 则
	      $f_n\weakstar f$ 在 $\sigma(E^\star,E)$ 中.
	\item 若 $f_n\weakstar f$ 在 $\sigma(E^\star,E)$ 中, 则 $(\norm{f_n})$ 有界,
	      且 $\norm{f}\le\liminf\norm{f_n}$.
	\item 若 $f_n\weakstar f$ 在 $\sigma(E^\star,E)$ 中, 且 $x_n\to x$ 在 $E$ 中强
	      收敛, 则 $\dual{f_n}{x_n}\to\dual{f}{x}$.
\end{enumerate}
\end{keyproposition}

\begin{proof}
照抄命题  \ref{prop:3.5} 的证明.
\end{proof}

\begin{remark}
\label{rem:3.9}
假设 $f_n\weakstar f$ 在 $\sigma(E^\star,E)$ 中 (甚至在
$\sigma(E^\star,E^{\star\star})$ 中弱收敛), 且 $x_n\weakto x$ 在
$\sigma(E,E^\star)$ 中.一般不能断定 $\dual{f_n}{x_n}\to\dual{f}{x}$ (在 Hilbert
空间中很容易构造出例子).
\end{remark}

\begin{remark}
\label{rem:3.10}
当 $E$ 为有限维空间时, $E^\star$ 上的三个拓扑 (强, 弱, \wstar{}) 重合.事实上,
典范单射 $J:E\to E^{\star\star}$ (见第  \ref{sec:1.3} 节) 是满射 (因为
$\dim E=\dim E^{\star\star}$), 因此 $\sigma(E^\star,E)=\sigma(E^\star,E^{\star\star})$.
\end{remark}

\begin{starproposition}
\label{prop:3.14}
设 $\varphi:E^\star\to\R$ 为关于\wstar{}拓扑连续的线性泛函.则存在 $x_0\in E$,
使
\begin{equation*}
\varphi(f)=\dual{f}{x_0}\quad\forall f\in E^\star.
\end{equation*}
\end{starproposition}

其证明依赖下述有用的代数引理:

\begin{lemma}
\label{lem:3.2}
设 $X$ 为向量空间, $\varphi,\varphi_1,\varphi_2,\dots,\varphi_k$ 为 $X$ 上的
$(k+1)$ 个线性泛函, 使得
\begin{equation}
[\varphi_i(v)=0\ \ \forall i=1,2,\dots,k]\ \Rightarrow\ [\varphi(v)=0].
\label{eq:3.2}
\end{equation}
则存在常数 $\lambda_1,\lambda_2,\dots,\lambda_k\in\R$, 使
$\varphi=\sum_{i=1}^k\lambda_i\varphi_i$.
\end{lemma}


\begin{proof}[引理 \ref{lem:3.2}]
考虑由
\begin{equation*}
F(u)=[\varphi(u),\varphi_1(u),\varphi_2(u),\dots,\varphi_k(u)]
\end{equation*}
定义的映射 $F:X\to\R^{k+1}$.由假设 \eqref{eq:3.2} 可知,
$a=[1,0,0,\dots,0]$ 不属于 $R(F)$.于是可用 $\R^{k+1}$ 中的某个超平面把
$\{a\}$ 与 $R(F)$ 严格分离; 即存在常数
$\lambda,\lambda_1,\lambda_2,\dots,\lambda_k$ 与 $\alpha$, 使
\begin{equation*}
\lambda<\alpha<\lambda\varphi(u)+\sum_{i=1}^k\lambda_i\varphi_i(u)
\quad\forall u\in X.
\end{equation*}
由此推出
\begin{equation*}
\lambda\varphi(u)+\sum_{i=1}^k\lambda_i\varphi_i(u)=0\quad\forall u\in X,
\end{equation*}
并且 $\lambda<0$ (从而 $\lambda\ne0$).
\end{proof}

\begin{proof}[命题 \ref{prop:3.14}]
由于 $\varphi$ 关于\wstar{}拓扑连续, 存在 $0$ 关于 $\sigma(E^\star,E)$ 的邻域
$V$, 使
\begin{equation*}
\abs{\varphi(f)}<1\quad\forall f\in V.
\end{equation*}
总可设
\begin{equation*}
V=\{f\in E^\star;\ \abs{\dual{f}{x_i}}<\varepsilon\ \ \forall i=1,2,\dots,k\},
\end{equation*}
其中 $x_i\in E$, $\varepsilon>0$.特别地,
\begin{equation*}
[\dual{f}{x_i}=0\ \ \forall i=1,2,\dots,k]\ \Rightarrow\ [\varphi(f)=0].
\end{equation*}
由引理  \ref{lem:3.2} 得
\begin{equation*}
\varphi(f)=\sum_{i=1}^k\lambda_i\dual{f}{x_i}
=\Bigl\langle f,\sum_{i=1}^k\lambda_ix_i\Bigr\rangle\quad\forall f\in E^\star.
\end{equation*}
\end{proof}

\begin{starcorollary}
\label{cor:3.15}
假设 $H$ 是 $E^\star$ 中的超平面, 且在 $\sigma(E^\star,E)$ 中闭.则 $H$ 具有形式
\begin{equation*}
H=\{f\in E^\star;\ \dual{f}{x_0}=\alpha\},
\end{equation*}
其中 $x_0\in E$, $x_0\ne0$, $\alpha\in\R$.
\end{starcorollary}

\begin{proof}
$H$ 可写成
\begin{equation*}
H=\{f\in E^\star;\ \varphi(f)=\alpha\},
\end{equation*}
其中 $\varphi$ 是 $E^\star$ 上的线性泛函, $\varphi\not\equiv0$.取
$f_0\notin H$, 并取 $f_0$ 关于拓扑 $\sigma(E^\star,E)$ 的邻域 $V$, 使
$V\subset H^c$.可设
\begin{equation*}
V=\{f\in E^\star;\ \abs{\dual{f-f_0}{x_i}}<\varepsilon\ \ \forall i=1,2,\dots,k\}.
\end{equation*}
由于 $V$ 凸, 我们有
\begin{equation}
\varphi(f)<\alpha\quad\forall f\in V
\label{eq:3.3}
\end{equation}
或
\begin{equation*}
\varphi(f)>\alpha\quad\forall f\in V.
\tag{3$'$}
\label{eq:3.3p}
\end{equation*}
例如设 \eqref{eq:3.3} 成立, 则得
\begin{equation*}
\varphi(g)<\alpha-\varphi(f_0)\quad\forall g\in W=V-f_0,
\end{equation*}
而由于 $-W=W$, 我们推出
\begin{equation}
\abs{\varphi(g)}\le\abs{\alpha-\varphi(f_0)}\quad\forall g\in W.
\label{eq:3.4}
\end{equation}
由 \eqref{eq:3.4} 可知 $\varphi$ 在 $0$ 处关于拓扑 $\sigma(E^\star,E)$ 连续 (因为
$W$ 是 $0$ 的邻域).应用命题  \ref{prop:3.14}, 我们断定存在 $x_0\in E$, 使
\begin{equation*}
\varphi(f)=\dual{f}{x_0}\quad\forall f\in E^\star.
\end{equation*}
\end{proof}

\begin{remark}
\label{rem:3.11}
假设典范单射 $J:E\to E^{\star\star}$ 不是满射.则拓扑 $\sigma(E^\star,E)$ 严格
粗于拓扑 $\sigma(E^\star,E^{\star\star})$.例如, 取 $\xi\in E^{\star\star}$,
$\xi\notin J(E)$.则集合
\begin{equation*}
H=\{f\in E^\star;\ \dual{\xi}{f}=0\}
\end{equation*}
在 $\sigma(E^\star,E^{\star\star})$ 中闭, 但由推论  \ref{cor:3.15} 可知它在
$\sigma(E^\star,E)$ 中不闭.从这个例子还可知: 在强拓扑中闭的凸集, 在\wstar{}
拓扑中不必闭.$E^\star$ 中的闭凸集有两类:
\begin{enumerate}[label=(\alph*), leftmargin=2.0em, itemsep=0.15em]
	\item 强闭的凸集 ($=$ 在拓扑 $\sigma(E^\star,E^{\star\star})$ 中闭, 由定理
	       \ref{thm:3.7});
	\item 在 $\sigma(E^\star,E)$ 中闭的凸集.
\end{enumerate}
\end{remark}

\index{BanachAlaogluBourbakidingli@Banach--Alaoglu--Bourbaki 定理 (Banach--Alaoglu--Bourbaki theorem)}
\begin{keytheorem}[Banach--Alaoglu--Bourbaki]
\label{thm:3.16}
闭单位球
\begin{equation*}
B_{E^\star}=\{f\in E^\star;\ \norm{f}\le1\}
\end{equation*}
在\wstar{}拓扑 $\sigma(E^\star,E)$ 中紧.
\end{keytheorem}

\begin{remark}
\label{rem:3.12}
$B_{E^\star}$ 的紧性是\wstar{}拓扑最本质的性质; 另见注记 \ref{rem:3.8}.
\end{remark}

\begin{proof}
考虑 Cartesian 积 $Y=\R^E$, 它由所有从 $E$ 到 $\R$ 的映射组成; 把 $Y$ 的元素记作
$\omega=(\omega_x)_{x\in E}$, 其中 $\omega_x\in\R$.空间 $Y$ 赋以标准\emph{乘积
拓扑} (见例如 H.~L.~Royden \bref{Royden}{1}, J.~R.~Munkres \bref{Munkres}{1}, A.~Knapp \bref{Knapp}{1} 或 J.~Dixmier \bref{Dixmier}{1}), 即 $Y$ 上与映射族
$\omega\mapsto\omega_x$ (当 $x$ 取遍 $E$ 时) 相伴的最粗拓扑; 这当然与逐点收敛
拓扑相同 (见例如 J.~R.~Munkres \bref{Munkres}{1}).以下总给 $E^\star$ 配上
\wstar{}拓扑 $\sigma(E^\star,E)$.

由于 $E^\star$ 由从 $E$ 到 $\R$ 的特殊映射 (即连续线性映射) 组成, 我们可以把
$E^\star$ 看作 $Y$ 的子集.更确切地说, 设 $\Phi:E^\star\to Y$ 为从 $E^\star$ 到
$Y$ 的典范单射, 使得 $\Phi(f)=(\omega_x)_{x\in E}$, 其中 $\omega_x=\dual{f}{x}$.
显然 $\Phi$ 是从 $E^\star$ 到 $Y$ 的连续映射 (用命题  \ref{prop:3.2}, 并注意对每个
固定的 $x\in E$, 映射 $f\in E^\star\mapsto(\Phi(f))_x=\dual{f}{x}$ 连续).逆映射
$\Phi^{-1}$ 也是从 $\Phi(E^\star)$ (赋以 $Y$ 的拓扑) 到 $E^\star$ 的连续映射:
事实上, 再次用命题  \ref{prop:3.2}, 只须验证对每个固定的 $x\in E$, 映射
$\omega\mapsto\dual{\Phi^{-1}(\omega)}{x}$ 在 $\Phi(E^\star)$ 上连续; 这是显然的,
因为 $\dual{\Phi^{-1}(\omega)}{x}=\omega_x$ (注意 $\omega=\Phi(f)$ 对某个
$f\in E^\star$ 成立, 且 $\dual{\Phi^{-1}(\omega)}{x}=\dual{f}{x}=\omega_x$).
换言之, $\Phi$ 是从 $E^\star$ 到 $\Phi(E^\star)$ 的同胚.另一方面, 显然
$\Phi(B_{E^\star})=K$, 其中 $K$ 由
\begin{equation*}
K=\Bigl\{\omega\in Y\ \Big|\ \abs{\omega_x}\le\norm{x},\
\omega_{x+y}=\omega_x+\omega_y\
\text{且}\ \omega_{\lambda x}=\lambda\omega_x\ \forall\lambda\in\R,\ \forall x,y\in E\Bigr\}
\end{equation*}
定义.为完成定理  \ref{thm:3.16} 的证明, 只须验证 $K$ 是 $Y$ 的紧子集.把 $K$
写成 $K=K_1\cap K_2$, 其中
\begin{equation*}
K_1=\{\omega\in Y;\ \abs{\omega_x}\le\norm{x}\ \ \forall x\in E\}
\end{equation*}
且
\begin{equation*}
K_2=\bigl\{\omega\in Y;\ \omega_{x+y}=\omega_x+\omega_y\
\text{且}\ \omega_{\lambda x}=\lambda\omega_x\ \forall\lambda\in\R,\ \forall x,y\in E\bigr\}.
\end{equation*}
集合 $K_1$ 也可写成紧区间的乘积
\begin{equation*}
K_1=\prod_{x\in E}[-\norm{x},+\norm{x}].
\end{equation*}
回忆: (任意) 紧空间的乘积紧——这是 Tychonoff 的一个深刻定理; 见例如
H.~L.~Royden \bref{Royden}{1}, G.~B.~Folland \bref{Folland}{2}, J.~R.~Munkres \bref{Munkres}{1}, A.~Knapp \bref{Knapp}{1} 或 J.~Dixmier
\bref{Dixmier}{1}.因此 $K_1$ 紧.另一方面, $K_2$ 在 $Y$ 中闭; 事实上, 对每个固定
的 $\lambda\in\R$, $x,y\in E$, 集合
\begin{equation*}
A_{x,y}=\{\omega\in Y;\ \omega_{x+y}-\omega_x-\omega_y=0\},
\end{equation*}
\begin{equation*}
B_{\lambda,x}=\{\omega\in Y;\ \omega_{\lambda x}-\lambda\omega_x=0\}
\end{equation*}
在 $Y$ 中闭 (因为映射 $\omega\mapsto\omega_{x+y}-\omega_x-\omega_y$ 与
$\omega\mapsto\omega_{\lambda x}-\lambda\omega_x$ 在 $Y$ 上连续), 而 $K_2$ 可写成
\begin{equation*}
K_2=\Bigl(\bigcap_{x,y\in E}A_{x,y}\Bigr)\cap
\Bigl(\bigcap_{\substack{x\in E\ \lambda\in\R}}B_{\lambda,x}\Bigr).
\end{equation*}
最后, $K$ 紧, 因为它是紧集 ($K_1$) 与闭集 ($K_2$) 的交.
\end{proof}
\section{自反空间}
\label{sec:3.5}

\begin{definition}
设 $E$ 为 Banach 空间, $J:E\to E^{\star\star}$ 为从 $E$ 到 $E^{\star\star}$ 的
典范单射 (见第  \ref{sec:1.3} 节).若 $J$ 是满射, 即 $J(E)=E^{\star\star}$, 则称
空间 $E$ 是\emph{自反的}.
\end{definition}

当 $E$ 自反时, 通常把 $E^{\star\star}$ 与 $E$ 等同起来.

\begin{remark}
\label{rem:3.13}
分析中许多重要的空间都是自反的.显然, 有限维空间自反 (因为
$\dim E=\dim E^\star=\dim E^{\star\star}$).我们将在第~
\ref{chap:4}~章 (另见第~
\ref{chap:11}~章) 看到,
当 $1<p<\infty$ 时 $L^p$ (以及 $\ell^p$) 空间自反.在第~
\ref{chap:5}~章我们将看到 Hilbert
空间自反.然而, 分析中同样重要的一些空间不是自反的; 例如:
\begin{itemize}[leftmargin=1.6em, itemsep=0.15em]
	\item $L^1$ 与 $L^\infty$ (以及 $\ell^1$, $\ell^\infty$) 不自反 (见第~
\ref{chap:4}~章与
	      第~
\ref{chap:11}~章);
	\item $C(K)$, 即无穷紧度量空间 $K$ 上的连续函数空间, 不自反 (见习题
	       \ref{exer:3.25}).
\end{itemize}
\end{remark}

\begin{starremark}
\label{rem:3.14}
在上述定义中使用 $J$ 是本质的.R.~C.~James \bref{James}{1} 构造了一个引人注目的
非自反空间的例子, 它具有下述性质: 存在从 $E$ 到 $E^{\star\star}$ 的满射等距.
\end{starremark}

我们的下一个结果描述自反空间的一个基本性质:

\begin{keytheorem}[Kakutani]
\label{thm:3.17}
设 $E$ 为 Banach 空间.则 $E$ 自反当且仅当
\begin{equation*}
B_E=\{x\in E;\ \norm{x}\le1\}
\end{equation*}
在弱拓扑 $\sigma(E,E^\star)$ 中紧.
\end{keytheorem}

\begin{proof}
先设 $E$ 自反, 于是 $J(B_E)=B_{E^{\star\star}}$.由定理  \ref{thm:3.16} 已知
$B_{E^{\star\star}}$ 在拓扑 $\sigma(E^{\star\star},E^\star)$ 中紧.因此只须验证
$J^{-1}$ 是从赋以 $\sigma(E^{\star\star},E^\star)$ 的 $E^{\star\star}$ 到赋以
$\sigma(E,E^\star)$ 的 $E$ 的连续映射.由命题  \ref{prop:3.2}, 我们只须证明对每个
固定的 $f\in E^\star$, 映射 $\xi\mapsto\dual{f}{J^{-1}\xi}$ 在赋以
$\sigma(E^{\star\star},E^\star)$ 的 $E^{\star\star}$ 上连续.但
$\dual{f}{J^{-1}\xi}=\dual{\xi}{f}$, 而映射 $\xi\mapsto\dual{\xi}{f}$ 确实在
$E^{\star\star}$ 上关于拓扑 $\sigma(E^{\star\star},E^\star)$ 连续.因此我们证明了
$B_E$ 在 $\sigma(E,E^\star)$ 中紧.

逆命题更为精细, 它依赖下述两个引理:
\end{proof}

\begin{lemma}[Helly]
\label{lem:3.3}
设 $E$ 为 Banach 空间.给定 $E^\star$ 中的 $f_1,f_2,\dots,f_k$ 以及 $\R$ 中的
$\gamma_1,\gamma_2,\dots,\gamma_k$.则下述性质等价:
\begin{enumerate}[label=(\roman*), leftmargin=2.4em, itemsep=0.15em]
	\item $\forall\varepsilon>0$, $\exists x_\varepsilon\in E$, 使
	      $\norm{x_\varepsilon}\le1$ 且
	      \begin{equation*}
	      \abs{\dual{f_i}{x_\varepsilon}-\gamma_i}<\varepsilon
	      \quad\forall i=1,2,\dots,k,
	      \end{equation*}
	\item $\abs{\sum_{i=1}^k\beta_i\gamma_i}
	      \le\norm{\sum_{i=1}^k\beta_i f_i}
	      \quad\forall\beta_1,\beta_2,\dots,\beta_k\in\R$.
\end{enumerate}
\end{lemma}

\begin{proof}
(i) $\Rightarrow$ (ii). 固定 $\beta_1,\beta_2,\dots,\beta_k\in\R$, 并置
$S=\sum_{i=1}^k\abs{\beta_i}$.由 (i) 得
\begin{equation*}
\abs{\sum_{i=1}^k\beta_i\dual{f_i}{x_\varepsilon}
-\sum_{i=1}^k\beta_i\gamma_i}\le\varepsilon S,
\end{equation*}
从而
\begin{equation*}
\abs{\sum_{i=1}^k\beta_i\gamma_i}
\le\norm{\sum_{i=1}^k\beta_i f_i}\norm{x_\varepsilon}+\varepsilon S
\le\norm{\sum_{i=1}^k\beta_i f_i}+\varepsilon S.
\end{equation*}
由于这对每个 $\varepsilon>0$ 都成立, 便得到 (ii).

(ii) $\Rightarrow$ (i). 置 $\gamma=(\gamma_1,\gamma_2,\dots,\gamma_k)\in\R^k$,
考虑由
\begin{equation*}
\varphi(x)=(\dual{f_1}{x},\dots,\dual{f_k}{x})
\end{equation*}
定义的映射 $\varphi:E\to\R^k$.性质 (i) 说的正是
$\gamma\in\overline{\varphi(B_E)}$.反设 (i) 不成立, 即
$\gamma\notin\overline{\varphi(B_E)}$.于是 $\{\gamma\}$ 与
$\overline{\varphi(B_E)}$ 可用 $\R^k$ 中某个超平面严格分离; 即存在
$\beta=(\beta_1,\beta_2,\dots,\beta_k)\in\R^k$ 与 $\alpha\in\R$, 使
\begin{equation*}
\beta\cdot\varphi(x)<\alpha<\beta\cdot\gamma\quad\forall x\in B_E.
\end{equation*}
由此推出
\begin{equation*}
\Bigl\langle\sum_{i=1}^k\beta_i f_i,x\Bigr\rangle<\alpha<\sum_{i=1}^k\beta_i\gamma_i
\quad\forall x\in B_E,
\end{equation*}
于是
\begin{equation*}
\norm{\sum_{i=1}^k\beta_i f_i}\le\alpha<\sum_{i=1}^k\beta_i\gamma_i,
\end{equation*}
这与 (ii) 矛盾.
\end{proof}

\begin{lemma}[Goldstine]
\label{lem:3.4}
设 $E$ 为任一 Banach 空间.则 $J(B_E)$ 在 $B_{E^{\star\star}}$ 中关于拓扑
$\sigma(E^{\star\star},E^\star)$ 稠密; 从而 $J(E)$ 在 $E^{\star\star}$ 中关于拓扑
$\sigma(E^{\star\star},E^\star)$ 稠密.
\end{lemma}

\begin{proof}
设 $\xi\in B_{E^{\star\star}}$, $V$ 为 $\xi$ 关于拓扑
$\sigma(E^{\star\star},E^\star)$ 的邻域.须证明
$V\cap J(B_E)\ne\emptyset$.如通常那样, 可设 $V$ 具有形式
\begin{equation*}
V=\bigl\{\eta\in E^{\star\star};\ \abs{\dual{\eta-\xi}{f_i}}<\varepsilon
\ \ \forall i=1,2,\dots,k\bigr\},
\end{equation*}
其中 $f_1,f_2,\dots,f_k$ 为 $E^\star$ 中给定的元素, $\varepsilon>0$.须找出
$x\in B_E$, 使 $J(x)\in V$, 即
\begin{equation*}
\abs{\dual{f_i}{x}-\dual{\xi}{f_i}}<\varepsilon\quad\forall i=1,2,\dots,k.
\end{equation*}
置 $\gamma_i=\dual{\xi}{f_i}$.由引理  \ref{lem:3.3}, 只须验证
\begin{equation*}
\abs{\sum_{i=1}^k\beta_i\gamma_i}\le\norm{\sum_{i=1}^k\beta_i f_i},
\end{equation*}
这是显然的, 因为
$\sum_{i=1}^k\beta_i\gamma_i=\bigl\langle\xi,\sum_{i=1}^k\beta_i f_i\bigr\rangle$
且 $\norm{\xi}\le1$.
\end{proof}

\begin{remark}
\label{rem:3.15}
注意 $J(B_E)$ 在 $B_{E^{\star\star}}$ 中关于强拓扑闭.事实上, 若
$\xi_n=J(x_n)\to\xi$, 则看出 $(x_n)$ 是 $B_E$ 中的 Cauchy 列 (因为 $J$ 是等距),
从而 $x_n\to x$, 于是 $\xi=Jx$.由此可知, 除非 $J(B_E)=B_{E^{\star\star}}$
(即 $E$ 自反), $J(B_E)$ 在 $B_{E^{\star\star}}$ 中关于强拓扑不稠密.
\end{remark}

\begin{remark}
\label{rem:3.16}
引理  \ref{lem:3.4} 的另一证明见问题 9 (它基于 $E^{\star\star}$ 中 Hahn--Banach
定理的一个变形).
\end{remark}

\begin{proof}[定理 \ref{thm:3.17} 的续证]
典范单射 $J:E\to E^{\star\star}$ 总是从 $\sigma(E,E^\star)$ 到
$\sigma(E^{\star\star},E^\star)$ 连续的, 因为对每个固定的 $f\in E^\star$, 映射
$x\mapsto\dual{Jx}{f}=\dual{f}{x}$ 关于 $\sigma(E,E^\star)$ 连续.假设 $B_E$ 在
拓扑 $\sigma(E,E^\star)$ 中紧, 我们推出 $J(B_E)$ 在 $E^{\star\star}$ 中关于拓扑
$\sigma(E^{\star\star},E^\star)$ 紧——从而闭.另一方面, 由引理  \ref{lem:3.4},
$J(B_E)$ 在同一拓扑下在 $B_{E^{\star\star}}$ 中稠密.于是
$J(B_E)=B_{E^{\star\star}}$, 从而 $J(E)=E^{\star\star}$.
\end{proof}

关于自反空间的紧性, 我们还有下述两个结果:

\begin{keytheorem}
\label{thm:3.18}
假设 $E$ 为自反 Banach 空间, $(x_n)$ 为 $E$ 中的有界序列.则存在子列
$(x_{n_k})$, 它在弱拓扑 $\sigma(E,E^\star)$ 中收敛.
\end{keytheorem}

逆命题也成立, 即:

\begin{startheorem}[Eberlein--\v{S}mulian]
\label{thm:3.19}
假设 $E$ 为 Banach 空间, 使 $E$ 中每个有界序列都有弱收敛 (在
$\sigma(E,E^\star)$ 中) 的子列.则 $E$ 自反.
\end{startheorem}

定理  \ref{thm:3.18} 的证明需要绕道可分空间, 将在第  \ref{sec:3.6} 节给出.定理
 \ref{thm:3.19} 的证明相当精细, 从略; 见例如 R.~Holmes \bref{Holmes}{1}, K.~Yosida \bref{Yosida}{1}, N.~Dunford--J.~T.~Schwartz \bref{DunfordSchwartz}{1}, J.~Diestel \bref{Diestel}{2} 或问题 10.

\begin{remark}
\label{rem:3.17}
为弄清定理  \ref{thm:3.17},  \ref{thm:3.18} 与  \ref{thm:3.19} 之间的联系, 回忆
下述事实是有用的:
\begin{enumerate}[label=(\roman*), leftmargin=2.4em, itemsep=0.2em]
	\item 若 $X$ 为度量空间, 则
	      \begin{equation*}
	      [X\ \text{紧}]\ \iff\ [X\ \text{中每个序列都有收敛子列}].
	      \end{equation*}
	\item 存在紧拓扑空间 $X$ 以及 $X$ 中没有任何收敛子列的序列.典型的例子是
	      $X=B_{E^\star}$, 它在拓扑 $\sigma(E^\star,E)$ 中紧; 当 $E=\ell^\infty$ 时
	      很容易构造 $X$ 中没有收敛子列的序列 (见习题  \ref{exer:3.18}).
	\item 若拓扑空间 $X$ 具有``每个序列都有收敛子列''的性质, 则 $X$ 不必紧.
\end{enumerate}
\end{remark}

下面是自反空间的另一些性质.

\begin{keyproposition}
\label{prop:3.20}
假设 $E$ 为自反 Banach 空间, $M\subset E$ 为 $E$ 的闭线性子空间.则 $M$ 自反.
\end{keyproposition}

\begin{proof}
空间 $M$ (赋以 $E$ 的范数) 先验地有两个不同的弱拓扑:
\begin{enumerate}[label=(\alph*), leftmargin=2.0em, itemsep=0.15em]
	\item 由 $\sigma(E,E^\star)$ 诱导的拓扑;
	\item 它自己的弱拓扑 $\sigma(M,M^\star)$.
\end{enumerate}
事实上, 这两个拓扑相同 (因为由 Hahn--Banach 定理, $M$ 上每个连续线性泛函都是
$E$ 上某个连续线性泛函在 $M$ 上的限制).由定理  \ref{thm:3.17}, 须验证 $B_M$ 在
拓扑 $\sigma(M,M^\star)$ 中紧, 或在拓扑 $\sigma(E,E^\star)$ 中紧.然而, $B_E$ 在
拓扑 $\sigma(E,E^\star)$ 中紧, 而 $M$ 在拓扑 $\sigma(E,E^\star)$ 中闭 (由定理
 \ref{thm:3.7}).因此 $B_M$ 在拓扑 $\sigma(E,E^\star)$ 中紧.
\end{proof}

\begin{corollary}
\label{cor:3.21}
Banach 空间 $E$ 自反当且仅当其共轭空间 $E^\star$ 自反.
\end{corollary}

\begin{proof}
$E$ 自反 $\Rightarrow$ $E^\star$ 自反.证明的想法很简单, 因为粗略地说, 我们有
$E^{\star\star}=E\Rightarrow E^{\star\star\star}=E^\star$.更确切地说, 设 $J$ 为
从 $E$ 到 $E^{\star\star}$ 的典范同构.给定 $\varphi\in E^{\star\star\star}$.
映射 $x\mapsto\dual{\varphi}{Jx}$ 是 $E$ 上的连续线性泛函.记它为 $f\in E^\star$,
于是
\begin{equation*}
\dual{\varphi}{Jx}=\dual{f}{x}\quad\forall x\in E.
\end{equation*}
但我们也有
\begin{equation*}
\dual{\varphi}{Jx}=\dual{Jx}{f}\quad\forall x\in E.
\end{equation*}
由于 $J$ 满射, 我们推出
\begin{equation*}
\dual{\varphi}{\xi}=\dual{\xi}{f}\quad\forall\xi\in E^{\star\star},
\end{equation*}
这正好说明从 $E^\star$ 到 $E^{\star\star\star}$ 的典范单射是满射.

$E^\star$ 自反 $\Rightarrow$ $E$ 自反.由上一段我们已知 $E^{\star\star}$ 自反.
由于 $J(E)$ 是 $E^{\star\star}$ 中强闭的子空间, 由命题  \ref{prop:3.20} 推得
$J(E)$ 自反.因此 $E$ 自反.\footnote{显然, 若 $E$ 与 $F$ 为 Banach 空间, 且 $T$
是从 $E$ 到 $F$ 的线性满射等距, 则 $E$ 自反当且仅当 $F$ 自反.当然, 这与注记
\ref{rem:3.14} 并不矛盾!}
\end{proof}

\begin{keycorollary}
\label{cor:3.22}
设 $E$ 为自反 Banach 空间, $K\subset E$ 为 $E$ 的有界闭凸子集.则 $K$ 在拓扑
$\sigma(E,E^\star)$ 中紧.
\end{keycorollary}

\begin{proof}
$K$ 关于拓扑 $\sigma(E,E^\star)$ 闭 (由定理  \ref{thm:3.7}).另一方面, 存在常数
$m$ 使 $K\subset mB_E$, 而 $mB_E$ 在 $\sigma(E,E^\star)$ 中紧 (由定理
 \ref{thm:3.17}).
\end{proof}

\begin{keycorollary}
\label{cor:3.23}
设 $E$ 为自反 Banach 空间, $A\subset E$ 为 $E$ 的非空闭凸子集.设
$\varphi:A\to(-\infty,+\infty]$ 为凸下半连续函数, 且 $\varphi\not\equiv+\infty$, 并
设
\begin{equation}
\lim_{\substack{x\in A\\ \norm{x}\to\infty}}\varphi(x)=+\infty
\label{eq:3.5}
\end{equation}
(若 $A$ 有界则无此假设).则 $\varphi$ 在 $A$ 上达到最小值, 即存在 $x_0\in A$, 使
\begin{equation*}
\varphi(x_0)=\min_A\varphi.
\end{equation*}
\end{keycorollary}

\begin{proof}
任取 $a\in A$ 使 $\varphi(a)<+\infty$, 并考虑集合
\begin{equation*}
\tilde{A}=\{x\in A;\ \varphi(x)\le\varphi(a)\}.
\end{equation*}
则 $\tilde{A}$ 闭, 凸且有界 (由 \eqref{eq:3.5}), 从而它在拓扑 $\sigma(E,E^\star)$
中紧 (由推论  \ref{cor:3.22}).另一方面, $\varphi$ 在拓扑 $\sigma(E,E^\star)$ 中也
下半连续 (由推论  \ref{cor:3.9}).于是 $\varphi$ 在 $\tilde{A}$ 上达到最小值 (见
第~
\ref{chap:1}~章下半连续函数定义之后的性质 5), 即存在 $x_0\in\tilde{A}$, 使
\begin{equation*}
\varphi(x_0)\le\varphi(x)\quad\forall x\in\tilde{A}.
\end{equation*}
若 $x\in A\setminus\tilde{A}$, 则 $\varphi(x_0)\le\varphi(a)<\varphi(x)$; 因此
\begin{equation*}
\varphi(x_0)\le\varphi(x)\quad\forall x\in A.
\end{equation*}
\end{proof}

\begin{remark}
\label{rem:3.18}
推论  \ref{cor:3.23} 是自反空间与凸函数在变分法与最优化中的许多问题里如此重要的
主要原因.
\end{remark}

\begin{theorem}
\label{thm:3.24}
设 $E$ 与 $F$ 为两个自反 Banach 空间.设 $A:D(A)\subset E\to F$ 为稠定且闭的
无界线性算子.则 $D(A^\star)$ 在 $F^\star$ 中稠密.于是 $A^{\star\star}$ 有定义
($A^{\star\star}:D(A^{\star\star})\subset E^{\star\star}\to F^{\star\star}$),
它也可看作从 $E$ 到 $F$ 的无界算子.并且有
\begin{equation*}
\boxed{A^{\star\star}=A.}
\end{equation*}
\end{theorem}

\begin{proof}
\emph{1. $D(A^\star)$ 在 $F^\star$ 中稠密.} 设 $\varphi$ 为 $F^\star$ 上的连续
线性泛函, 它在 $D(A^\star)$ 上为零.由推论  \ref{cor:1.8}, 只须证明 $\varphi\equiv0$
在 $F^\star$ 上.由于 $F$ 自反, $\varphi\in F$, 且有
\begin{equation}
\dual{w}{\varphi}=0\quad\forall w\in D(A^\star).
\label{eq:3.6}
\end{equation}
若 $\varphi\ne0$, 则 $[0,\varphi]\notin G(A)$ (在 $E\times F$ 中).于是可用
$E\times F$ 中的闭超平面把 $[0,\varphi]$ 与 $G(A)$ 严格分离; 即存在
$[f,v]\in E^\star\times F^\star$ 与 $\alpha\in\R$, 使
\begin{equation*}
\dual{f}{u}+\dual{v}{Au}<\alpha<\dual{v}{\varphi}\quad\forall u\in D(A).
\end{equation*}
由此推出
\begin{equation*}
\dual{f}{u}+\dual{v}{Au}=0\quad\forall u\in D(A)
\end{equation*}
且
\begin{equation*}
\dual{v}{\varphi}\ne0.
\end{equation*}
于是 $v\in D(A^\star)$, 而在 \eqref{eq:3.6} 中取 $w=v$ 便导致矛盾.

\emph{2. $A^{\star\star}=A$.} 回忆 (见第  \ref{sec:2.6} 节)
\begin{equation*}
I[G(A^\star)]=G(A)^\perp
\end{equation*}
且
\begin{equation*}
I[G(A^{\star\star})]=G(A^\star)^\perp.
\end{equation*}
由此
\begin{equation*}
G(A^{\star\star})=G(A)^{\perp\perp}=G(A),
\end{equation*}
因为 $A$ 闭.
\end{proof}

\section{可分空间}
\label{sec:3.6}

\begin{definition}
若度量空间 $E$ 存在可数稠密子集 $D\subset E$, 则称 $E$ 是\emph{可分的}.
\end{definition}

分析中许多重要的空间都是可分的.显然, 有限维空间可分.我们将在第~
\ref{chap:4}~章 (另见
第~
\ref{chap:11}~章) 看到, 当 $1\le p<\infty$ 时 $L^p$ (以及 $\ell^p$) 空间可分.此外,
$C(K)$, 即紧度量空间 $K$ 上的连续函数空间, 也可分 (见问题 24).然而,
$L^\infty$ 与 $\ell^\infty$ 不可分 (见第~
\ref{chap:4}~章与第~
\ref{chap:11}~章).

\begin{proposition}
\label{prop:3.25}
设 $E$ 为可分度量空间, $F\subset E$ 为任一子集.则 $F$ 也可分.
\end{proposition}

\begin{proof}
设 $(u_n)$ 为 $E$ 的可数稠密子集.设 $(r_m)$ 为任一趋于 $0$ 的正数列.只要
$B(u_n,r_m)\cap F$ 非空, 就任取一点 $a_{m,n}\in B(u_n,r_m)\cap F$.集合
$(a_{m,n})$ 可数且在 $F$ 中稠密.
\end{proof}

\begin{theorem}
\label{thm:3.26}
设 $E$ 为 Banach 空间, 且 $E^\star$ 可分.则 $E$ 可分.
\end{theorem}

\begin{remark}
\label{rem:3.19}
逆命题不成立.我们将在第~
\ref{chap:4}~章看到, $E=L^1$ 可分, 但其共轭空间
$E^\star=L^\infty$ 不可分.
\end{remark}

\begin{proof}
设 $(f_n)_{n\ge1}$ 在 $E^\star$ 中可数稠密.由于
\begin{equation*}
\norm{f_n}=\sup_{\substack{x\in E\\ \norm{x}\le1}}\dual{f_n}{x},
\end{equation*}
我们可以找到 $x_n\in E$, 使
\begin{equation*}
\norm{x_n}=1\quad\text{且}\quad\dual{f_n}{x_n}\ge\frac12\norm{f_n}.
\end{equation*}
以 $L_0$ 记 $(x_n)_{n\ge1}$ 在 $\Q$ 上张成的向量空间; 即 $L_0$ 由 $(x_n)_{n\ge1}$
的, 系数在 $\Q$ 中的所有有限线性组合组成.我们断言 $L_0$ 可数.事实上, 对每个
整数 $n$, 设 $\Phi_n$ 为 $(x_k)_{1\le k\le n}$ 在 $\Q$ 上张成的向量空间.显然
$\Phi_n$ 可数, 且 $L_0=\bigcup_{n\ge1}\Phi_n$.

以 $L$ 记 $(x_n)_{n\ge1}$ 在 $\R$ 上张成的向量空间.当然, $L_0$ 是 $L$ 的稠密
子集.我们断言 $L$ 是 $E$ 的稠密子空间——这就完成了证明 ($L_0$ 将是 $E$ 的可数
稠密子集).设 $f\in E^\star$ 为在 $L$ 上为零的连续线性泛函; 由推论
 \ref{cor:1.8}, 须证明 $f=0$.给定任意 $\varepsilon>0$, 存在整数 $N$, 使
$\norm{f-f_N}<\varepsilon$.我们有
\begin{equation*}
\frac12\norm{f_N}\le\dual{f_N}{x_N}=\dual{f_N-f}{x_N}<\varepsilon
\end{equation*}
(因为 $\dual{f}{x_N}=0$).于是 $\norm{f}\le\norm{f-f_N}+\norm{f_N}<3\varepsilon$.
因此 $f=0$.
\end{proof}

\begin{corollary}
\label{cor:3.27}
设 $E$ 为 Banach 空间.则
\begin{equation*}
[E\ \text{自反且可分}]\ \iff\ [E^\star\ \text{自反且可分}].
\end{equation*}
\end{corollary}

\begin{proof}
由推论  \ref{cor:3.21} 与定理  \ref{thm:3.26} 已知
\begin{equation*}
[E^\star\ \text{自反且可分}]\ \Rightarrow\ [E\ \text{自反且可分}].
\end{equation*}
反之, 若 $E$ 自反且可分, 则 $E^{\star\star}=J(E)$ 也自反且可分; 因此 $E^\star$
自反且可分.
\end{proof}

可分性与弱拓扑的可度量化密切相关.回忆: 若拓扑空间 $X$ 上存在一个度量诱导 $X$
的拓扑, 则称 $X$ 是\emph{可度量化的}.

\begin{theorem}
\label{thm:3.28}
设 $E$ 为可分 Banach 空间.则 $B_{E^\star}$ 在\wstar{}拓扑
$\sigma(E^\star,E)$ 中可度量化.

反之, 若 $B_{E^\star}$ 在 $\sigma(E^\star,E)$ 中可度量化, 则 $E$ 可分.
\end{theorem}

还有一个``对偶''的命题.

\begin{theorem}
\label{thm:3.29}
设 $E$ 为 Banach 空间, 且 $E^\star$ 可分.则 $B_E$ 在弱拓扑
$\sigma(E,E^\star)$ 中可度量化.

反之, 若 $B_E$ 在 $\sigma(E,E^\star)$ 中可度量化, 则 $E^\star$ 可分.
\end{theorem}

\begin{proof}[定理 \ref{thm:3.28}]
设 $(x_n)_{n\ge1}$ 为 $B_E$ 的可数稠密子集.对每个 $f\in E^\star$ 置
\begin{equation*}
[f]=\sum_{n=1}^\infty\frac{1}{2^n}\abs{\dual{f}{x_n}}.
\end{equation*}
显然, $[\cdot]$ 是 $E^\star$ 上的范数, 且 $[f]\le\norm{f}$.令
$d(f,g)=[f-g]$ 为相应的度量.我们将证明: $d$ 在 $B_{E^\star}$ 上诱导的拓扑与
$\sigma(E^\star,E)$ 限制在 $B_{E^\star}$ 上的拓扑相同.

(a) 设 $f_0\in B_{E^\star}$, $V$ 为 $f_0$ 关于 $\sigma(E^\star,E)$ 的邻域.须找出
$r>0$, 使
\begin{equation*}
U=\{f\in B_{E^\star};\ d(f,f_0)<r\}\subset V.
\end{equation*}
如通常那样, 可设 $V$ 具有形式
\begin{equation*}
V=\{f\in B_{E^\star};\ \abs{\dual{f-f_0}{y_i}}<\varepsilon
\ \ \forall i=1,2,\dots,k\},
\end{equation*}
其中 $\varepsilon>0$, $y_1,y_2,\dots,y_k\in E$.不妨设
$\norm{y_i}\le1$ 对每个 $i=1,2,\dots,k$ 成立.对每个 $i$ 存在整数 $n_i$, 使
\begin{equation*}
\norm{y_i-x_{n_i}}<\varepsilon/4
\end{equation*}
(因为集合 $(x_n)_{n\ge1}$ 在 $B_E$ 中稠密).取 $r>0$ 充分小, 使
\begin{equation*}
2^{n_i}r<\varepsilon/2\quad\forall i=1,2,\dots,k.
\end{equation*}
我们断言: 对这样的 $r$, $U\subset V$.事实上, 若 $d(f,f_0)<r$, 则
\begin{equation*}
\frac{1}{2^{n_i}}\abs{\dual{f-f_0}{x_{n_i}}}<r\quad\forall i=1,2,\dots,k,
\end{equation*}
从而对每个 $i=1,2,\dots,k$,
\begin{equation*}
\abs{\dual{f-f_0}{y_i}}
=\abs{\dual{f-f_0}{y_i-x_{n_i}}+\dual{f-f_0}{x_{n_i}}}
<\frac{\varepsilon}{2}+\frac{\varepsilon}{2}.
\end{equation*}
由此 $f\in V$.

(b) 设 $f_0\in B_{E^\star}$.给定 $r>0$, 须找出 $f_0$ 关于 $\sigma(E^\star,E)$ 的
邻域 $V$, 使
\begin{equation*}
V\subset U=\{f\in B_{E^\star};\ d(f,f_0)<r\}.
\end{equation*}
我们取 $V$ 为
\begin{equation*}
V=\{f\in B_{E^\star};\ \abs{\dual{f-f_0}{x_i}}<\varepsilon
\ \ \forall i=1,2,\dots,k\},
\end{equation*}
其中 $\varepsilon$ 与 $k$ 待定, 使得 $V\subset U$.对 $f\in V$ 有
\begin{equation*}
\begin{aligned}
d(f,f_0)&=\sum_{n=1}^k\frac{1}{2^n}\abs{\dual{f-f_0}{x_n}}
+\sum_{n=k+1}^\infty\frac{1}{2^n}\abs{\dual{f-f_0}{x_n}}\\
&<\varepsilon+2\sum_{n=k+1}^\infty\frac{1}{2^n}
=\varepsilon+\frac{1}{2^{k-1}}.
\end{aligned}
\end{equation*}
因此只须取 $\varepsilon=\frac r2$, 并取 $k$ 充分大, 使
$\frac{1}{2^{k-1}}<\frac r2$.
\end{proof}

\begin{proof}[定理 \ref{thm:3.28} 的逆命题]
$\star$ 反过来, 设 $B_{E^\star}$ 在 $\sigma(E^\star,E)$ 中可度量化, 我们来证明
$E$ 可分.置
\begin{equation*}
U_n=\{f\in B_{E^\star};\ d(f,0)<1/n\},
\end{equation*}
并设 $V_n$ 为 $0$ 在 $\sigma(E^\star,E)$ 中的邻域, 使 $V_n\subset U_n$.可设 $V_n$
具有形式
\begin{equation*}
V_n=\{f\in B_{E^\star};\ \abs{\dual{f}{x}}<\varepsilon_n\ \ \forall x\in\Phi_n\},
\end{equation*}
其中 $\varepsilon_n>0$, $\Phi_n$ 为 $E$ 的有限子集.置
\begin{equation*}
D=\bigcup_{n=1}^\infty\Phi_n,
\end{equation*}
则 $D$ 可数.

我们断言: 由 $D$ 张成的向量空间在 $E$ 中稠密 (由此推出 $E$ 可分).事实上, 设
$f\in E^\star$ 满足 $\dual{f}{x}=0$, $\forall x\in D$.则 $f\in V_n$,
$\forall n$, 从而 $f\in U_n$, $\forall n$, 于是 $f=0$.
\end{proof}

\begin{proof}[定理 \ref{thm:3.29}]
蕴含关系
\begin{equation*}
[E^\star\ \text{可分}]\ \Rightarrow\ [B_E\ \text{在}\ \sigma(E,E^\star)\ \text{中可度量化}]
\end{equation*}
的证明与上面完全相同——只需把 $E$ 与 $E^\star$ 的角色互换.逆命题的证明更为
精细 (请找出上面的证明在哪里失效); 我们引用 N.~Dunford--J.~T.~Schwartz
\bref{DunfordSchwartz}{1} 或习题  \ref{exer:3.24}.
\end{proof}

\begin{remark}
\label{rem:3.20}
应当再次强调 (见注记 \ref{rem:3.3}): 在无穷维空间中, 整个 $E$ (相应地,
$E^\star$) 上的弱拓扑 $\sigma(E,E^\star)$ (相应地, \wstar{}拓扑
$\sigma(E^\star,E)$) 不可度量化; 见习题  \ref{exer:3.8}.特别地, 范数
$[\cdot]$ 在整个 $E^\star$ 上诱导的拓扑与\wstar{}拓扑并不相同.
\end{remark}

\begin{corollary}
\label{cor:3.30}
设 $E$ 为可分 Banach 空间, $(f_n)$ 为 $E^\star$ 中的有界序列.则存在子列
$(f_{n_k})$, 它在\wstar{}拓扑 $\sigma(E^\star,E)$ 中收敛.
\end{corollary}

\begin{proof}
不妨设 $\norm{f_n}\le1$ 对一切 $n$ 成立.集合 $B_{E^\star}$ 关于拓扑
$\sigma(E^\star,E)$ 紧且可度量化 (由定理  \ref{thm:3.16} 与  \ref{thm:3.28}).
结论由此得出.
\end{proof}

现在可以回到定理  \ref{thm:3.18} 的证明了:

\begin{proof}[定理 \ref{thm:3.18}]
设 $M_0$ 为由 $x_n$ 张成的向量空间, $M=\overline{M_0}$.显然 $M$ 可分 (见定理
 \ref{thm:3.26} 的证明).此外, $M$ 自反 (由命题  \ref{prop:3.20}).于是 $B_M$ 在
弱拓扑 $\sigma(M,M^\star)$ 中紧且可度量化, 因为 $M^\star$ 可分 (这里用到推论
 \ref{cor:3.27} 与定理  \ref{thm:3.29}).因此可以找到子列 $(x_{n_k})$, 它在
$\sigma(M,M^\star)$ 中弱收敛, 从而 $(x_{n_k})$ 也在 $\sigma(E,E^\star)$ 中弱收敛
(如命题  \ref{prop:3.20} 的证明中那样).
\end{proof}

\section{一致凸空间}
\label{sec:3.7}

\begin{definition}
若 Banach 空间满足
\begin{equation*}
\forall\varepsilon>0\ \exists\delta>0\ \text{使得}\
\Bigl[x,y\in E,\ \norm{x}\le1,\ \norm{y}\le1\ \text{且}\ \norm{x-y}>\varepsilon
\Bigr]\ \Rightarrow\ \Bigl\|\frac{x+y}{2}\Bigr\|<1-\delta,
\end{equation*}
则称它是\emph{一致凸的}.
\end{definition}

一致凸性是单位球的一个\emph{几何}性质: 若把一根长为 $\varepsilon>0$ 的直尺在单位
球中滑动, 则它的中点必须落在某个半径为 $(1-\delta)$ 的球内, 其中 $\delta>0$.特别
地, 单位\emph{球面}必须是``圆的'', 不能包含任何线段.

\begin{example}
\label{exa:3.3}
设 $E=\R^2$.范数 $\norm{x}_2=\bigl[\abs{x_1}^2+\abs{x_2}^2\bigr]^{1/2}$ 一致凸,
而范数 $\norm{x}_1=\abs{x_1}+\abs{x_2}$ 与
$\norm{x}_\infty=\max(\abs{x_1},\abs{x_2})$ 不是一致凸的.只要盯着单位球看,
如图  \ref{fig:3} 所示, 这一点就一目了然.
\end{example}

\begin{figure}[htbp]
	\centering
	\begin{tikzpicture}[scale=1.05]
		% 左图: l2 范数的单位圆
		\begin{scope}
			\draw[-{Stealth[length=2mm]}] (-1.55,0) -- (1.75,0);
			\draw[-{Stealth[length=2mm]}] (0,-1.55) -- (0,1.75);
			\draw[thick] (0,0) circle (1.0);
			\fill (0,0) circle (1.5pt);
			\node[below=6mm, align=center] at (0,-1.0) {\small 关于范数 $\norm{\cdot}_2$ 的 $E$ 的单位球};
		\end{scope}
		% 右图: l1 范数的单位菱形
		\begin{scope}[xshift=6.0cm]
			\draw[-{Stealth[length=2mm]}] (-1.55,0) -- (1.75,0);
			\draw[-{Stealth[length=2mm]}] (0,-1.55) -- (0,1.75);
			\draw[thick] (1,0) -- (0,1) -- (-1,0) -- (0,-1) -- cycle;
			\fill (0,0) circle (1.5pt);
			\fill (1,0) circle (1.5pt);
			\fill (-1,0) circle (1.5pt);
			\fill (0,1) circle (1.5pt);
			\fill (0,-1) circle (1.5pt);
			\node[below=6mm, align=center] at (0,-1.0) {\small 关于范数 $\norm{\cdot}_1$ 的 $E$ 的单位球};
		\end{scope}
	\end{tikzpicture}
	\caption{}
	\label{fig:3}
\end{figure}

\begin{example}
\label{exa:3.4}
我们将在第~
\ref{chap:4}~章与第~
\ref{chap:5}~章看到, 当 $1<p<\infty$ 时 $L^p$ 空间一致凸, 而 Hilbert
空间也一致凸.
\end{example}

\index{MilmanPettisdingli@Milman--Pettis 定理 (Milman--Pettis theorem)}
\begin{keytheorem}[Milman--Pettis]
\label{thm:3.31}
每个一致凸 Banach 空间都自反.
\end{keytheorem}

\begin{remark}
\label{rem:3.21}
一致凸性是范数的一个\emph{几何}性质; 等价的范数不必一致凸.另一方面, 自反性是
一个\emph{拓扑}性质: 自反空间关于等价范数仍然自反.定理  \ref{thm:3.31} 的一个
引人注目的特点是: 一个几何性质蕴含了一个拓扑性质.一致凸性常被用作证明自反性的
工具; 但它不是终极工具——存在一些古怪的自反空间, 它们不允许任何一致凸的等价
范数!
\end{remark}

\begin{proof}
设 $\xi\in E^{\star\star}$ 且 $\norm{\xi}=1$.须证明 $\xi\in J(B_E)$.由于
$J(B_E)$ 在 $E^{\star\star}$ 中关于强拓扑闭, 只须证明
\begin{equation}
\forall\varepsilon>0\ \exists x\in B_E\ \text{使得}\ \norm{\xi-J(x)}\le\varepsilon.
\label{eq:3.7}
\end{equation}
固定 $\varepsilon>0$, 设 $\delta>0$ 为一致凸性的模.选取某个 $f\in E^\star$, 使
$\norm{f}=1$ 且
\begin{equation}
\dual{\xi}{f}>1-(\delta/2)
\label{eq:3.8}
\end{equation}
(这是可能的, 因为 $\norm{\xi}=1$).置
\begin{equation*}
V=\{\eta\in E^{\star\star};\ \abs{\dual{\eta-\xi}{f}}<\delta/2\},
\end{equation*}
于是 $V$ 是 $\xi$ 在拓扑 $\sigma(E^{\star\star},E^\star)$ 中的一个邻域.由于
$J(B_E)$ 关于 $\sigma(E^{\star\star},E^\star)$ 在 $B_{E^{\star\star}}$ 中稠密
(引理  \ref{lem:3.4}), 我们知道 $V\cap J(B_E)\ne\emptyset$, 从而存在 $x\in B_E$
使 $J(x)\in V$.我们断言这个 $x$ 满足 \eqref{eq:3.7}.反设
$\norm{\xi-Jx}>\varepsilon$, 即 $\xi\in(Jx+\varepsilon B_{E^{\star\star}})^c=W$.
集合 $W$ 也是 $\xi$ 在拓扑 $\sigma(E^{\star\star},E^\star)$ 中的一个邻域 (因为
$B_{E^{\star\star}}$ 在 $\sigma(E^{\star\star},E^\star)$ 中闭).再次利用引理
 \ref{lem:3.4}, 我们知道 $V\cap W\cap J(B_E)\ne\phi$, 即存在 $y\in B_E$ 使
$J(y)\in V\cap W$.由 $J(x),J(y)\in V$, 得
\begin{equation*}
\abs{\dual{f}{x}-\dual{\xi}{f}}<\delta/2
\end{equation*}
以及
\begin{equation*}
\abs{\dual{f}{y}-\dual{\xi}{f}}<\delta/2.
\end{equation*}
把这两个不等式相加, 得到
\begin{equation*}
2\dual{\xi}{f}<\dual{f}{x+y}+\delta\le\norm{x+y}+\delta.
\end{equation*}
结合 \eqref{eq:3.8}, 得到
\begin{equation*}
\Bigl\|\frac{x+y}{2}\Bigr\|>1-\delta.
\end{equation*}
由一致凸性推出 $\norm{x-y}\le\varepsilon$; 这是荒谬的, 因为 $J(y)\in W$
(即 $\norm{x-y}>\varepsilon$).
\end{proof}

最后给出一致凸空间的一个有用性质.

\begin{proposition}
\index{RadonRieszxingzhi@Radon--Riesz 性质 (Radon--Riesz property)}
\label{prop:3.32}
假设 $E$ 是一致凸 Banach 空间.设 $(x_n)$ 为 $E$ 中的序列, 使
$x_n\weakto x$ 弱收敛于 $\sigma(E,E^\star)$, 且
\begin{equation*}
\limsup\norm{x_n}\le\norm{x}.
\end{equation*}
则 $x_n\to x$ 强收敛.
\end{proposition}

\begin{proof}
总可设 $x\ne0$ (否则结论显然).置
\begin{equation*}
\lambda_n=\max(\norm{x_n},\norm{x}),\quad
y_n=\lambda_n^{-1}x_n,\quad y=\norm{x}^{-1}x,
\end{equation*}
于是 $\lambda_n\to\norm{x}$ 且 $y_n\weakto y$ 弱收敛于 $\sigma(E,E^\star)$.
由此
\begin{equation*}
\norm{y}\le\liminf\norm{(y_n+y)/2}
\end{equation*}
(见命题  \ref{prop:3.5}(iii)).另一方面, $\norm{y}=1$ 且 $\norm{y_n}\le1$, 因此
事实上 $\norm{(y_n+y)/2}\to1$.由一致凸性推出 $\norm{y_n-y}\to0$, 从而
$x_n\to x$ 强收敛.
\end{proof}

\section*{第 3 章评注}
\addcontentsline{toc}{section}{第 3 章评注}
\markright{第 3 章评注}

\cmtnum{1.}
拓扑 $\sigma(E,E^\star)$, $\sigma(E^\star,E)$ 等都是局部凸拓扑.作为这样的拓扑,
它们享有局部凸空间的一切性质; 例如 Hahn--Banach 定理 (几何形式), Krein--Milman 定理等仍然成立; 参见例如 N.~Bourbaki \bref{Bourbaki}{1}, A.~Knapp \bref{Knapp}{1}, 以及问题 9.

\cmtnum{2.}
下面是\wstar{}拓扑的另一个值得提及的引人注目的性质.

\begin{startheorem}[Banach--Dieudonn\'e--Krein--\v{S}mulian]
\label{thm:3.33}
设 $E$ 为 Banach 空间, $C\subset E^\star$ 凸.假设对每个整数 $n$, 集合
$C\cap(nB_{E^\star})$ 关于拓扑 $\sigma(E^\star,E)$ 闭.则 $C$ 关于拓扑
$\sigma(E^\star,E)$ 闭.
\end{startheorem}

其证明可在例如 N.~Bourbaki \bref{Bourbaki}{1}, R.~Larsen \bref{Larsen}{1}, R.~Holmes \bref{Holmes}{1}, N.~Dunford--J.~T.~Schwartz \bref{DunfordSchwartz}{1}, H.~Schaefer \bref{Schaefer}{1} 以及问题 11 中找到.上述文献也包含了大量与
Eberlein--\v{S}mulian 定理 (定理  \ref{thm:3.19}) 相关的材料.

\cmtnum{3.}
\emph{对偶向量空间}理论——它推广了对偶 $\dual{E}{E^\star}$——在四十年代末与
五十年代初非常流行, 尤其是在与广义函数理论的联系中.若在 $X\times Y$ 上存在一个
分离点的双线性形式 $\langle\,,\rangle$ (即 $\forall x\ne0\ \exists y$ 使
$\dual{x}{y}\ne0$, 且 $\forall y\ne0\ \exists x$ 使 $\dual{x}{y}\ne0$), 则称
两个向量空间 $X$ 与 $Y$ 处于对偶之中.在 $X$ (或 $Y$) 上可以定义许多拓扑, 例如
弱拓扑 $\sigma(X,Y)$, Mackey 拓扑 $\tau(X,Y)$ 以及强拓扑 $\beta(X,Y)$.这些
拓扑在非 Banach 空间——例如广义函数理论中使用的空间——里很有意义.关于这一
主题, 读者可参阅例如 N.~Bourbaki \bref{Bourbaki}{1}, H.~Schaefer
\bref{Schaefer}{1}, G.~K\"othe \bref{Kothe}{1}, F.~Treves \bref{Treves}{1}, J.~Kelley--I.~Namioka \bref{KelleyNamioka}{1}, R.~Edwards \bref{Edwards}{1}, J.~Horv\'ath \bref{Horvath}{1} 等.

\cmtnum{4.}
可分性, 自反性与一致凸性这些性质也与函数 $x\mapsto\norm{x}$ 的可微性性质有关
(参见例如 J.~Diestel \bref{Diestel}{1}, B.~Beauzamy \bref{Beauzamy}{1} 以及
问题 13).关于具有良好几何性质的等价范数的存在性, 已有大量研究.例如, 怎样判断
一个 Banach 空间是否允许一个等价的\emph{一致凸}范数? 这一信息有多大用处? (这样的
空间称为\emph{超自反的}; 参见例如 J.~Diestel \bref{Diestel}{1} 或 B.~Beauzamy
\bref{Beauzamy}{1}.)\ Banach 空间的几何自六十年代初以来蓬勃发展, 已成为一个活跃
的领域, 与 A.~Dvoretzky, A.~Grothendieck, R.~C.~James, J.~Lindenstrauss, V.~Milman, L.~Tzafriri (以及他们在以色列的学派), A.~Pelczynski, P.~Enflo, L.~Schwartz (以及他包括 G.~Pisier, B.~Maurey, B.~Beauzamy 在内的学派), W.~B.~Johnson, H.~P.~Rosenthal, J.~Bourgain, D.~Preiss, M.~Talagrand, T.~Gowers 等许多人的名字联系在一起.关于这一主题, 读者可参阅 B.~Beauzamy
\bref{Beauzamy}{1}, J.~Diestel \bref{Diestel}{1}, \bref{Diestel}{2}, J.~Lindenstrauss--L.~Tzafriri \bref{LindenstraussTzafriri}{2}, L.~Schwartz
\bref{Schwartz}{2}, R.~Deville--G.~Godefroy--V.~Zizler \bref{Deville}{1}, Y.~Benyamini--J.~Lindenstrauss \bref{Benyamini}{1}, F.~Albiac--N.~Kalton
\bref{Albiac}{1}, A.~Pietsch \bref{Pietsch}{1} 等著作.

\section*{第 3 章习题}
\addcontentsline{toc}{section}{第 3 章习题}
\markright{第 3 章习题}

\begin{exercise}
\label{exer:3.1}
设 $E$ 为 Banach 空间, $A\subset E$ 为在弱拓扑 $\sigma(E,E^\star)$ 中紧的子集.
证明 $A$ 有界.
\end{exercise}

\begin{exercise}
\label{exer:3.2}
设 $E$ 为 Banach 空间, $(x_n)$ 为 $E$ 中的序列, 使 $x_n\weakto x$ 在弱拓扑
$\sigma(E,E^\star)$ 中弱收敛.置
\begin{equation*}
\sigma_n=\frac1n(x_1+x_2+\cdots+x_n).
\end{equation*}
证明 $\sigma_n\weakto x$ 在弱拓扑 $\sigma(E,E^\star)$ 中弱收敛.
\end{exercise}

\begin{exercise}
\label{exer:3.3}
设 $E$ 为 Banach 空间, $A\subset E$ 为凸子集.证明 $A$ 在强拓扑中的闭包与它在
弱拓扑 $\sigma(E,E^\star)$ 中的闭包相同.
\end{exercise}

\begin{exercise}
\label{exer:3.4}
设 $E$ 为 Banach 空间, $(x_n)$ 为 $E$ 中的序列, 使 $x_n\weakto x$ 在弱拓扑
$\sigma(E,E^\star)$ 中弱收敛.

\begin{enumerate}[label=\arabic*., leftmargin=2.2em, itemsep=0.25em]
	\item 证明存在 $E$ 中的序列 $(y_n)$, 使得
	      \begin{equation*}
	      \text{(a)}\quad y_n\in\operatorname{conv}\Bigl(\bigcup_{i=n}^\infty\{x_i\}\Bigr)
	      \quad\forall n
	      \end{equation*}
	      并且
	      \begin{equation*}
	      \text{(b)}\quad y_n\to x\quad\text{强收敛}.
	      \end{equation*}
	\item 证明存在 $E$ 中的序列 $(z_n)$, 使得
	      \begin{equation*}
	      \text{(a$'$)}\quad z_n\in\operatorname{conv}\Bigl(\bigcup_{i=1}^n\{x_i\}\Bigr)
	      \quad\forall n
	      \end{equation*}
	      并且
	      \begin{equation*}
	      \text{(b$'$)}\quad z_n\to x\quad\text{强收敛}.
	      \end{equation*}
\end{enumerate}
\end{exercise}

\begin{exercise}
\label{exer:3.5}
设 $E$ 为 Banach 空间, $K\subset E$ 为在强拓扑中紧的子集.设 $(x_n)$ 为 $K$ 中的
序列, 使 $x_n\weakto x$ 弱收敛于 $\sigma(E,E^\star)$.证明 $x_n\to x$ 强收敛.

[提示: 用反证法.]
\end{exercise}

\begin{exercise}
\label{exer:3.6}
设 $X$ 为拓扑空间, $E$ 为 Banach 空间.设 $u,v:X\to E$ 为两个连续映射, 其中
$E$ 配上弱拓扑 $\sigma(E,E^\star)$.

\begin{enumerate}[label=\arabic*., leftmargin=2.2em, itemsep=0.25em]
	\item 证明映射 $x\mapsto u(x)+v(x)$ 是从 $X$ 到配上 $\sigma(E,E^\star)$ 的
	      $E$ 的连续映射.
	\item 设 $a:X\to\R$ 为连续函数.证明映射 $x\mapsto a(x)u(x)$ 是从 $X$ 到
	      配上 $\sigma(E,E^\star)$ 的 $E$ 的连续映射.
\end{enumerate}
\end{exercise}

\begin{exercise}
\label{exer:3.7}
设 $E$ 为 Banach 空间, $A\subset E$ 为在弱拓扑 $\sigma(E,E^\star)$ 中闭的子集.
设 $B\subset E$ 为在弱拓扑 $\sigma(E,E^\star)$ 中紧的子集.

\begin{enumerate}[label=\arabic*., leftmargin=2.2em, itemsep=0.25em]
	\item 证明 $A+B$ 在 $\sigma(E,E^\star)$ 中闭.
	\item 再假设 $A$ 与 $B$ 凸, 非空且不相交.证明存在一个闭超平面严格分离 $A$
	      与 $B$.
\end{enumerate}
\end{exercise}

\begin{exercise}
\label{exer:3.8}
设 $E$ 为无穷维 Banach 空间.我们的目的是证明配上弱拓扑的 $E$ 不可度量化.
反设 $E$ 上存在一个度量 $d(x,y)$, 它在 $E$ 上诱导出与 $\sigma(E,E^\star)$ 相同的
拓扑.

\begin{enumerate}[label=\arabic*., leftmargin=2.2em, itemsep=0.25em]
	\item 对每个整数 $k\ge1$, 用 $V_k$ 记 $0$ 在拓扑 $\sigma(E,E^\star)$ 中的
	      一个邻域, 使得
	      \begin{equation*}
	      V_k\subset\Bigl\{x\in E;\ d(x,0)<\frac1k\Bigr\}.
	      \end{equation*}
	      证明存在 $E^\star$ 中的序列 $(f_n)$, 使得每个 $g\in E^\star$ 都是
	      这些 $f_n$ 的 (有限) 线性组合.

	      [提示: 利用引理  \ref{lem:3.2}.]
	\item 推出 $E^\star$ 是有限维的.

	      [提示: 像习题 1.5 那样利用 Baire 纲定理.]
	\item 得出结论.
	\item 用类似的方法证明配上\wstar{}拓扑 $\sigma(E^\star,E)$ 的 $E^\star$
	      不可度量化.
\end{enumerate}
\end{exercise}

\begin{exercise}
\label{exer:3.9}
设 $E$ 为 Banach 空间, $M\subset E$ 为线性子空间, $f_0\in E^\star$.证明存在
某个 $g_0\in M^\perp$, 使得
\begin{equation*}
\inf_{g\in M^\perp}\norm{f_0-g}=\norm{f_0-g_0}.
\end{equation*}
建议两种方法:

\begin{enumerate}[label=\arabic*., leftmargin=2.2em, itemsep=0.25em]
	\item 利用定理  \ref{thm:1.12}.
	\item 利用\wstar{}拓扑 $\sigma(E^\star,E)$.
\end{enumerate}
\end{exercise}

\begin{exercise}
\label{exer:3.10}
设 $E$ 与 $F$ 为两个 Banach 空间.设 $T\in\LL(E,F)$, 从而
$T^\star\in\LL(F^\star,E^\star)$.证明 $T^\star$ 是从配上 $\sigma(F^\star,F)$
的 $F^\star$ 到配上 $\sigma(E^\star,E)$ 的 $E^\star$ 的连续映射.
\end{exercise}

\begin{exercise}
\label{exer:3.11}
设 $E$ 为 Banach 空间, $A:E\to E^\star$ 为定义在 $D(A)=E$ 上的单调映射;
见习题 2.6.假设对每个 $x,y\in E$, 映射
\begin{equation*}
t\in\R\mapsto\dual{A(x+ty)}{y}
\end{equation*}
在 $t=0$ 处连续.证明 $A$ 是从强拓扑下的 $E$ 到配上 $\sigma(E^\star,E)$ 的
$E^\star$ 的连续映射.
\end{exercise}

\begin{exercise}
\label{exer:3.12}
设 $E$ 为 Banach 空间, $x_0\in E$.设 $\varphi:E\to(-\infty,+\infty]$ 为凸的
l.s.c. 函数, $\varphi\not\equiv+\infty$.

\begin{enumerate}[label=\arabic*., leftmargin=2.2em, itemsep=0.25em]
	\item 证明下列性质等价:
	      \begin{equation*}
	      \text{(A)}\quad
	      \begin{aligned}
	      &\exists R,\ \exists M<+\infty\ \text{使得}\ \varphi(x)\le M,\\
	      &\forall x\in E\ \text{满足}\ \norm{x-x_0}\le R,
	      \end{aligned}
	      \end{equation*}
	      \begin{equation*}
	      \text{(B)}\quad
	      \lim_{\substack{f\in E^\star\\ \norm{f}\to\infty}}
	      \{\varphi^\star(f)-\dual{f}{x_0}\}=+\infty.
	      \end{equation*}
	\item 假设 (A) 或 (B) 成立, 证明
	      \begin{equation*}
	      \inf_{f\in E^\star}\{\varphi^\star(f)-\dual{f}{x_0}\}
	      \end{equation*}
	      可以达到.

	      [提示: 利用\wstar{}拓扑 $\sigma(E^\star,E)$ 或定理  \ref{thm:1.12}.]

	      这个 $\inf$ 的值是多少?
\end{enumerate}
\end{exercise}

\begin{exercise}
\label{exer:3.13}
设 $E$ 为 Banach 空间.设 $(x_n)$ 为 $E$ 中的序列, $x\in E$.置
\begin{equation*}
K_n=\overline{\operatorname{conv}\Bigl(\bigcup_{i=n}^\infty\{x_i\}\Bigr)}.
\end{equation*}

\begin{enumerate}[label=\arabic*., leftmargin=2.2em, itemsep=0.25em]
	\item 证明: 若 $x_n\weakto x$ 弱收敛于 $\sigma(E,E^\star)$, 则
	      \begin{equation*}
	      \bigcap_{n=1}^\infty K_n=\{x\}.
	      \end{equation*}
	\item 假设 $E$ 自反.证明: 若 $(x_n)$ 有界且
	      $\bigcap_{n=1}^\infty K_n=\{x\}$, 则 $x_n\weakto x$ 弱收敛于
	      $\sigma(E,E^\star)$.
	\item 假设 $E$ 有限维且 $\bigcap_{n=1}^\infty K_n=\{x\}$.证明
	      $x_n\to x$.

	      [注意这里并没有假设 $(x_n)$ 有界.]
	\item 在 $\ell^p$ ($1<p<\infty$, 见第~
\ref{chap:11}~章) 中构造一个序列 $(x_n)$, 使得
	      $\bigcap_{n=1}^\infty K_n=\{x\}$, 而 $(x_n)$ \emph{不}有界.

	      [第 3 问与第 4 问的结果得自 Guy Amram 与 Daniel Baffet.]
\end{enumerate}
\end{exercise}

\begin{exercise}
\label{exer:3.14}
设 $E$ 为自反 Banach 空间, $I$ 为指标集.考虑 $E^\star$ 中的一族 $(f_i)_{i\in I}$
与 $\R$ 中的一族 $(\alpha_i)_{i\in I}$.设 $M>0$.证明下列性质等价:
\begin{equation*}
\text{(A)}\quad
\begin{cases}
\text{存在某个}\ x\in E,\ \norm{x}\le M,\ \text{使得}\ \dual{f_i}{x}=\alpha_i\\
\text{对每个}\ i\in I.
\end{cases}
\end{equation*}
\begin{equation*}
\text{(B)}\quad
\begin{cases}
\text{对}\ \R\ \text{中任一满足}\ J\subset I,\ J\ \text{有限的族}\
(\beta_i)_{i\in J},\ \text{都有}\\
\abs{\sum_{i\in J}\beta_i\alpha_i}\le M\bigl\|\sum_{i\in J}\beta_if_i\bigr\|.
\end{cases}
\end{equation*}
与习题 1.10, 1.11 及引理  \ref{lem:3.3} 比较.
\end{exercise}

\begin{exercise}[凸集上测度的重心.]
\label{exer:3.15}
设 $E$ 为自反 Banach 空间, $K\subset E$ 有界, 闭且凸.下文给 $K$ 配上
$\sigma(E,E^\star)$, 于是 $K$ 紧.设 $F=C(K)$ 配上通常的范数.固定某个
$\mu\in F^\star$ 且 $\norm{\mu}=1$, 并假设 $\mu\ge0$, 即
\begin{equation*}
\dual{\mu}{u}\ge0\quad\forall u\in C(K),\ u\ge0\ \text{在}\ K\ \text{上}.
\end{equation*}
证明存在唯一的元素 $x_0\in K$, 使得
\begin{equation}
\dual{\mu}{f|_K}=\dual{f}{x_0}\quad\forall f\in E^\star.
\tag{1}
\end{equation}
[提示: 先找出某个满足 (1) 的 $x_0\in E$, 然后借助 Hahn--Banach 定理证明
$x_0\in K$.]
\end{exercise}

\begin{exercise}
\label{exer:3.16}
设 $E$ 为 Banach 空间.

\begin{enumerate}[label=\arabic*., leftmargin=2.2em, itemsep=0.25em]
	\item 设 $(f_n)$ 为 $E^\star$ 中的序列, 使得对每个 $x\in E$,
	      $\dual{f_n}{x}$ 收敛到一个极限.证明存在某个 $f\in E^\star$, 使得
	      $f_n\weakstar f$ 在 $\sigma(E^\star,E)$ 中\wstar{}收敛.
	\item 这里假设 $E$ 自反.设 $(x_n)$ 为 $E$ 中的序列, 使得对每个
	      $f\in E^\star$, $\dual{f}{x_n}$ 收敛到一个极限.证明存在某个 $x\in E$,
	      使得 $x_n\weakto x$ 在 $\sigma(E,E^\star)$ 中弱收敛.
	\item 在非自反空间 $E$ 中构造一个例子, 使第 2 问的结论失效.

	      [提示: 取 $E=c_0$ (见第~\ref{sec:11.3}~节) 以及
	      $x_n=(1,1,\dots,1_{(n)},0,0,\dots)$.]
\end{enumerate}
\end{exercise}

\begin{exercise}
\label{exer:3.17}
\begin{enumerate}[label=\arabic*., leftmargin=2.2em, itemsep=0.25em]
	\item 设 $(x_n)$ 为 $\ell^p$ 中的序列, $1\le p\le\infty$.假设 $x_n\weakto x$
	      在 $\sigma(\ell^p,\ell^{p'})$ 中弱收敛, 证明:
	      \begin{enumerate}[label=(\alph*), leftmargin=2.0em, itemsep=0.15em]
	      	\item $(x_n)$ 在 $\ell^p$ 中有界,
	      	\item 对每个 $i$ 有 $x^n_i\to x_i$ ($n\to\infty$), 其中
	      	      $x_n=(x^n_1,x^n_2,\dots,x^n_i,\dots)$,
	      	      $x=(x_1,x_2,\dots,x_i,\dots)$.
	      \end{enumerate}
	\item 反过来, 设 $(x_n)$ 为 $\ell^p$ 中的序列, $1<p\le\infty$.假设 (a) 与
	      (b) 成立 (对某个记作 $x_i$ 的极限).证明 $x\in\ell^p$ 且 $x_n\weakto x$
	      在 $\sigma(\ell^p,\ell^{p'})$ 中弱收敛.
\end{enumerate}
\end{exercise}

\begin{exercise}
\label{exer:3.18}
对每个整数 $n\ge1$, 置
\begin{equation*}
e_n=(0,0,\dots,1_{(n)},0,\dots).
\end{equation*}
\begin{enumerate}[label=\arabic*., leftmargin=2.2em, itemsep=0.25em]
	\item 证明 $e_n\weakto0$ 在弱拓扑 $\sigma(\ell^p,\ell^{p'})$ 中弱收敛,
	      $1<p\le\infty$.
	\item 证明不存在子列 $(e_{n_k})$ 在 $\ell^1$ 中关于 $\sigma(\ell^1,\ell^\infty)$
	      收敛.
	\item 构造一个 Banach 空间 $E$ 与 $E^\star$ 中的序列 $(f_n)$ 的例子, 使得
	      $\norm{f_n}=1\ \forall n$, 且 $(f_n)$ 没有在 $\sigma(E^\star,E)$ 中收敛的
	      子列.这与 $B_{E^\star}$ 在拓扑 $\sigma(E^\star,E)$ 中的紧性矛盾吗?

	      [提示: 取 $E=\ell^\infty$.]
\end{enumerate}
\end{exercise}

\begin{exercise}
\label{exer:3.19}
设 $E=\ell^p$ 且 $F=\ell^q$, 其中 $1<p<\infty$, $1<q<\infty$.设
$a:\R\to\R$ 为连续函数, 使得
\begin{equation*}
\abs{a(t)}\le C\abs{t}^{p/q}\quad\forall t\in\R.
\end{equation*}
给定
\begin{equation*}
x=(x_1,x_2,\dots,x_i,\dots)\in\ell^p,
\end{equation*}
置
\begin{equation*}
Ax=\bigl(a(x_1),a(x_2),\dots,a(x_i),\dots\bigr).
\end{equation*}
\begin{enumerate}[label=\arabic*., leftmargin=2.2em, itemsep=0.25em]
	\item 证明 $Ax\in\ell^q$, 且映射 $x\mapsto Ax$ 是从 $\ell^p$ (强) 到
	      $\ell^q$ (强) 的连续映射.
	\item 证明: 若 $(x^n)$ 为 $\ell^p$ 中的序列, 使 $x^n\weakto x$ 在
	      $\sigma(\ell^p,\ell^{p'})$ 中弱收敛, 则 $Ax^n\weakto Ax$ 在
	      $\sigma(\ell^q,\ell^{q'})$ 中弱收敛.
	\item 推出 $A$ 是从配上 $\sigma(E,E^\star)$ 的 $B_E$ 到配上
	      $\sigma(F,F^\star)$ 的 $F$ 的连续映射.
\end{enumerate}
\end{exercise}

\begin{exercise}
\label{exer:3.20}
设 $E$ 为 Banach 空间.

\begin{enumerate}[label=\arabic*., leftmargin=2.2em, itemsep=0.25em]
	\item 证明存在紧拓扑空间 $K$ 以及从 $E$ 到配上通常范数的 $C(K)$ 的等距
	      同构.

	      [提示: 取 $K=B_{E^\star}$ 配上 $\sigma(E^\star,E)$.]
	\item 假设 $E$ 可分, 证明存在从 $E$ 到 $\ell^\infty$ 的等距同构.
\end{enumerate}
\end{exercise}

\begin{exercise}
\label{exer:3.21}
设 $E$ 为可分 Banach 空间, $(f_n)$ 为 $E^\star$ 中的有界序列.直接证明——不利用
$E^\star$ 的可度量化性——存在一个子列 $(f_{n_k})$ 在 $\sigma(E^\star,E)$ 中
收敛.

[提示: 利用对角线过程.]
\end{exercise}

\begin{exercise}
\label{exer:3.22}
设 $E$ 为无穷维 Banach 空间, 满足下列假设之\emph{一}:
\begin{enumerate}[label=(\alph*), leftmargin=2.0em, itemsep=0.15em]
	\item $E^\star$ 可分,
	\item $E$ 自反.
\end{enumerate}
证明存在 $E$ 中的序列 $(x_n)$, 使得
\begin{equation*}
\norm{x_n}=1\quad\forall n
\qquad\text{且}\qquad
x_n\weakto0\ \text{弱收敛于}\ \sigma(E,E^\star).
\end{equation*}
\end{exercise}

\begin{exercise}
\label{exer:3.23}
若 $E$ 自反, 定理  \ref{thm:2.16} 的证明会变得容易得多.特别地, 找出
(b) $\Rightarrow$ (a) 的一个简单证明.
\end{exercise}

\begin{exercise}
\label{exer:3.24}
本习题的目的是勾勒定理  \ref{thm:3.29} 证明的一部分, 即: 若 Banach 空间 $E$
使得 $B_E$ 关于 $\sigma(E,E^\star)$ 可度量化, 则 $E^\star$ 可分.设 $d(x,y)$
为 $B_E$ 上的度量, 它在 $B_E$ 上诱导出与 $\sigma(E,E^\star)$ 相同的拓扑.置
\begin{equation*}
U_n=\Bigl\{x\in B_E;\ d(x,0)<\frac1n\Bigr\}.
\end{equation*}
设 $V_n$ 为 $0$ 关于 $\sigma(E,E^\star)$ 的邻域, 使 $V_n\subset U_n$.不妨设
$V_n$ 具有形式
\begin{equation*}
V_n=\{x\in E;\ \abs{\dual{f}{x}}<\varepsilon_n\quad\forall f\in\Phi_n\},
\end{equation*}
其中 $\varepsilon_n>0$, $\Phi_n\subset E^\star$ 是某个有限子集.置
$D=\bigcup_{n=1}^\infty\Phi_n$, 并设 $F$ 为 $D$ 张成的向量空间.我们断言
$F$ 在 $E^\star$ 中关于强拓扑稠密.反设 $\overline{F}\ne E^\star$.

\begin{enumerate}[label=\arabic*., leftmargin=2.2em, itemsep=0.25em]
	\item 证明存在某个 $\xi\in E^{\star\star}$ 与某个 $f_0\in E^\star$, 使得
	      \begin{equation*}
	      \dual{\xi}{f_0}>1,\qquad
	      \dual{\xi}{f}=0\quad\forall f\in F,\qquad
	      \text{且}\qquad \norm{\xi}=1.
	      \end{equation*}
	\item 置
	      \begin{equation*}
	      W=\Bigl\{x\in B_E;\ \abs{\dual{f_0}{x}}<\frac12\Bigr\}.
	      \end{equation*}
	      证明存在某个整数 $n_0\ge1$, 使 $V_{n_0}\subset W$.
	\item 证明存在 $x_1\in B_E$, 使得
	      \begin{equation*}
	      \begin{cases}
	      \abs{\dual{f}{x_1}-\dual{\xi}{f}}<\varepsilon_{n_0}
	      &\forall f\in\Phi_{n_0},\\[2pt]
	      \abs{\dual{f_0}{x_1}-\dual{\xi}{f_0}}<\dfrac12.
	      \end{cases}
	      \end{equation*}
	\item 推出 $x_1\in V_{n_0}$ 且 $\dual{f_0}{x_1}>\frac12$.
	\item 得出结论.
\end{enumerate}
\end{exercise}

\begin{exercise}
\label{exer:3.25}
设 $K$ 为紧度量空间, 且不是有限集.证明 $C(K)$ 不自反.

[提示: 设 $(a_n)$ 为 $K$ 中的序列, 使 $a_n\to a$ 且 $a_n\ne a\ \forall n$.
考虑线性泛函 $f(u)=\sum_{n=1}^\infty\frac1{2^n}u(a_n)$, $u\in C(K)$, 并按
习题 1.3 与 1.4 的方式论证.]
\end{exercise}

\begin{exercise}
\label{exer:3.26}
设 $F$ 为可分 Banach 空间, $(a_n)$ 为 $B_F$ 的稠密子集.考虑线性算子
$T:\ell^1\to F$,
\begin{equation*}
Tx=\sum_{i=1}^\infty x_ia_i
\qquad\text{其中}\quad x=(x_1,x_2,\dots,x_n,\dots)\in\ell^1.
\end{equation*}

\begin{enumerate}[label=\arabic*., leftmargin=2.2em, itemsep=0.25em]
	\item 证明 $T$ 有界且满射.
\end{enumerate}
下文再假设 $F$ 是无穷维的, 且 $F^\star$ 可分.

\begin{enumerate}[label=\arabic*., leftmargin=2.2em, itemsep=0.25em, start=2]
	\item 证明 $T$ 没有右逆.

	      [提示: 利用习题  \ref{exer:3.22} 与问题 8 的结果.]
	\item 推出 $N(T)$ 在 $\ell^1$ 中没有补空间.
	\item 确定 $E^\star$.
\end{enumerate}
\end{exercise}

\begin{exercise}
\label{exer:3.27}
设 $E$ 为可分 Banach 空间, 范数为 $\norm{\cdot}$.$E^\star$ 上的对偶范数也记作
$\norm{\cdot}$.本习题的目的是在 $E$ 上构造一个等价范数, 它严格凸, 且其对偶范数
也严格凸.

设 $(a_n)\subset B_E$ 为 $B_E$ 关于强拓扑的稠密子集.设
$(b_n)\subset B_{E^\star}$ 为 $B_{E^\star}$ 的可数子集, 它在 $B_{E^\star}$ 中
关于\wstar{}拓扑 $\sigma(E^\star,E)$ 稠密.这样的集合为何存在?

给定 $f\in E^\star$, 置
\begin{equation*}
\|f\|_1=\Bigl\{\|f\|^2+\sum_{n=1}^\infty\frac{1}{2^n}\abs{\dual{f}{a_n}}^2
\Bigr\}^{1/2}.
\end{equation*}
\begin{enumerate}[label=\arabic*., leftmargin=2.2em, itemsep=0.25em]
	\item 证明 $\|\cdot\|_1$ 是与 $\|\cdot\|$ 等价的范数.
	\item 证明 $\|\cdot\|_1$ 严格凸.

	      [提示: 利用习题 1.26.]
\end{enumerate}
给定 $x\in E$, 置
\begin{equation*}
\|x\|_2=\Bigl\{\|x\|_1^2+\sum_{n=1}^\infty\frac{1}{2^n}
\abs{\dual{b_n}{x}}^2\Bigr\}^{1/2},
\end{equation*}
其中 $\|x\|_1=\sup_{\|f\|_1\le1}\dual{f}{x}$.

\begin{enumerate}[label=\arabic*., leftmargin=2.2em, itemsep=0.25em, start=3]
	\item 证明 $\|\cdot\|_2$ 是与 $\|\cdot\|$ 等价的严格凸范数.
	\item 证明 $\|\cdot\|_2$ 的对偶范数也严格凸.

	      [提示: 利用习题 1.23 第 3 问的结果.]
	\item 基于问题 4 的结果找出另一种方法.
\end{enumerate}
\end{exercise}

\begin{exercise}
\label{exer:3.28}
设 $E$ 为一致凸 Banach 空间.用 $F$ 记从 $E$ 到 $E^\star$ 的 (多值) 对偶映射,
见推论  \ref{cor:1.3} 之后的注记 2 以及习题 1.1.

证明: 对每个 $f\in E^\star$, 存在唯一的 $x\in E$, 使得 $f\in Fx$.
\end{exercise}

\begin{exercise}
\label{exer:3.29}
设 $E$ 为一致凸 Banach 空间.

\begin{enumerate}[label=\arabic*., leftmargin=2.2em, itemsep=0.25em]
	\item 证明 $\forall M>0$, $\forall\varepsilon>0$, $\exists\delta>0$, 使得
	      \begin{equation*}
	      \Bigl\|\frac{x+y}{2}\Bigr\|^2\le\frac12\norm{x}^2+\frac12\norm{y}^2-\delta
	      \end{equation*}
	      \begin{equation*}
	      \forall x,y\in E\ \text{满足}\ \norm{x}\le M,\ \norm{y}\le M
	      \ \text{且}\ \norm{x-y}>\varepsilon.
	      \end{equation*}
	      [提示: 用反证法.]
	\item 当 $\|\cdot\|^2$ 换成 $\|\cdot\|^p$ ($1<p<\infty$) 时, 问同样的问题.
\end{enumerate}
\end{exercise}

\begin{exercise}
\label{exer:3.30}
设 $E$ 为范数 $\norm{\cdot}$ 的 Banach 空间.假设 $E$ 上存在一个等价的范数,
记作 $\abs{\cdot}$, 它一致凸.

证明: 对任意 $k>1$, 存在 $E$ 上的一致凸范数 $|\!|\!|\cdot|\!|\!|$, 使得
\begin{equation*}
\norm{x}\le|\!|\!|x|\!|\!|\le k\norm{x}\quad\forall x\in E.
\end{equation*}
[提示: 置 $|\!|\!|x|\!|\!|^2=\norm{x}^2+\alpha\abs{x}^2$, 其中 $\alpha>0$ 足够
小, 并利用习题  \ref{exer:3.29}.]

例: $E=\R^n$.
\end{exercise}

\begin{exercise}
\label{exer:3.31}
设 $E$ 为一致凸 Banach 空间.

\begin{enumerate}[label=\arabic*., leftmargin=2.2em, itemsep=0.25em]
	\item 证明
	      \begin{equation*}
	      \forall\varepsilon>0,\quad
	      \forall\alpha\in\Bigl(0,\frac12\Bigr),\quad
	      \exists\delta>0
	      \end{equation*}
	      使得
	      \begin{equation*}
	      \norm{tx+(1-t)y}\le1-\delta
	      \end{equation*}
	      \begin{equation*}
	      \forall t\in[\alpha,1-\alpha],\quad\forall x,y\in E
	      \ \text{满足}\ \norm{x}\le1,\ \norm{y}\le1\ \text{且}\
	      \norm{x-y}\ge\varepsilon.
	      \end{equation*}
	      [提示: 若 $\alpha\le t\le\frac12$, 则写
	      $tx+(1-t)y=\frac12(y+z)$.]
	\item 推出 $E$ 严格凸.
\end{enumerate}
\end{exercise}

\begin{exercise}[一致凸 Banach 空间中闭凸集上的投影.]
\label{exer:3.32}
设 $E$ 为一致凸 Banach 空间, $C\subset E$ 为非空闭凸集.

\begin{enumerate}[label=\arabic*., leftmargin=2.2em, itemsep=0.25em]
	\item 证明对每个 $x\in E$,
	      \begin{equation*}
	      \inf_{y\in C}\norm{x-y}
	      \end{equation*}
	      由 $C$ 中某个唯一点达到, 记作 $P_Cx$.
	\item 证明 $C$ 中每个极小化序列 $(y_n)$ 都强收敛到 $P_Cx$.
	\item 证明映射 $x\mapsto P_Cx$ 是从强拓扑下的 $E$ 到强拓扑下的 $E$ 的连续
	      映射.
	\item 更精确地, 证明 $P_C$ 在 $E$ 的有界子集上一致连续.

	      [提示: 利用习题  \ref{exer:3.29}.]
\end{enumerate}
设 $\varphi:E\to(-\infty,+\infty]$ 为凸的 l.s.c. 函数,
$\varphi\not\equiv+\infty$.

\begin{enumerate}[label=\arabic*., leftmargin=2.2em, itemsep=0.25em, start=5]
	\item 证明对每个 $x\in E$ 与每个整数 $n\ge1$,
	      \begin{equation*}
	      \inf_{y\in E}\bigl\{n\norm{x-y}^2+\varphi(y)\bigr\}
	      \end{equation*}
	      由某个唯一点达到, 记作 $y_n$.
	\item 证明 $y_n\to P_Cx$ ($n\to\infty$), 其中
	      $C=\overline{D(\varphi)}$.
\end{enumerate}
\end{exercise}
